% ============================================================================
% BOLUM 7: KAYNAK KOD REKONSTRUKSIYONU
% BOLUM 8: GUVENLIK ANALIZ ALGORITMALARI
% BOLUM 9: PERFORMANS VE OLCEKLENDIRME
% ============================================================================
% Dosya: chapters/ch07_08_09_recon_security_perf.tex
% Bagimlilklar: main.tex'teki tum paketler (tikz, amsmath, listings, tcolorbox)
% ============================================================================


% ############################################################################
%
%                    BOLUM 7 -- KAYNAK KOD REKONSTRUKSIYONU
%
% ############################################################################

\section{Kaynak Kod Rekonstruksiyonuna Giris}
\label{sec:recon-giris}

Tersine muhendisligin nihai hedefi, derlenmis veya obfuscate edilmis koddan
orijinal kaynak koda \emph{mumkun oldugunca yakin} bir temsil elde etmektir.
Karadul bu hedefe iki paralel pipeline ile ulasir:

\begin{enumerate}
  \item \textbf{Binary Reconstruction Pipeline:} Mach-O/ELF/PE binary'lerden
        Ghidra decompile ciktisi uzerinde C kodu rekonstruksiyonu.
  \item \textbf{JavaScript Reconstruction Pipeline:} Webpack/Rollup bundle'larindan
        minified JS kodun orijinal proje yapisina donusturulmesi.
\end{enumerate}

Her iki pipeline'da da ortak zorluk \emph{isim kurtarma} (name recovery)
problemidir: derleyici veya minifier tarafindan silinen orijinal degisken,
fonksiyon ve tip isimlerinin yeniden olusturulmasi.

\begin{tanim}[Rekonstruksiyon Kalitesi]
Bir rekonstruksiyon $\hat{S}$'nin kalitesi, orijinal kaynak $S$'e olan
yapissal ve semantik yakinligi ile olculur:
\begin{equation}
  Q(\hat{S}, S) = \alpha \cdot \text{StructSim}(\hat{S}, S)
                + \beta \cdot \text{NameSim}(\hat{S}, S)
                + \gamma \cdot \text{TypeSim}(\hat{S}, S)
  \label{eq:recon-quality}
\end{equation}
burada $\alpha + \beta + \gamma = 1$ agirliklari,
$\text{StructSim}$ kontrol akisi grafi (CFG) izomorfizm benzerligini,
$\text{NameSim}$ isimlendirme dogru\-lugunu (precision/recall),
$\text{TypeSim}$ ise tip kurtarma basar\-isini olcer.
\end{tanim}


% ============================================================================
\section{6 Katmanli C Isimlendirme Sistemi (c\_namer.py)}
\label{sec:c-namer}

Ghidra decompiler, fonksiyonlara \code{FUN\_00abcdef}, parametrelere
\code{param\_1}, yerel degiskenlere \code{local\_28}, global verilere
\code{DAT\_00123456} gibi otomatik isimler atar. Bu isimler anlasilabilirlik
acisindan neredeyse degersidir. \file{c\_namer.py} modulu, bu isimleri
altI katmanli bir strateji zinciriyle anlamli isimlere donusturur.

\subsection{Mimari Genel Bakis}

Sistem, her katmanin bilgi entropisi azaltma katkisiyla bir \emph{cascade
classifier} (kademeli siniflandirici) gibi calisir. Her katman bir oncekinin
bulamadigi isimleri bulmaya calisir ve her katmanin urettigi isimler farkli
guven seviyelerine sahiptir.

\begin{figure}[htbp]
\centering
\begin{tikzpicture}[
    layer/.style={
      draw=sectioncyan, fill=boxbg, text=white,
      minimum width=12cm, minimum height=1.2cm,
      rounded corners=3pt, font=\small\bfseries
    },
    arrow/.style={-{Stealth[length=5pt]}, thick, white},
    conf/.style={font=\scriptsize\ttfamily, text=green!70!white},
    note/.style={font=\scriptsize, text=gray!70!white, align=left},
    >=Stealth
]
  % Layers
  \node[layer] (L1) at (0,0) {Katman 1: Symbol / Export Table};
  \node[layer] (L2) at (0,-1.8) {Katman 2: String Context Analysis};
  \node[layer] (L3) at (0,-3.6) {Katman 3: API Call Pattern Matching};
  \node[layer] (L4) at (0,-5.4) {Katman 4: Call Graph Structure};
  \node[layer] (L5) at (0,-7.2) {Katman 5: Dataflow Analysis};
  \node[layer] (L6) at (0,-9.0) {Katman 6: Type-Based Inference};

  % Confidence annotations
  \node[conf, right] at (6.5,0) {conf: 0.95};
  \node[conf, right] at (6.5,-1.8) {conf: 0.70--0.90};
  \node[conf, right] at (6.5,-3.6) {conf: 0.60--0.80};
  \node[conf, right] at (6.5,-5.4) {conf: 0.50--0.70};
  \node[conf, right] at (6.5,-7.2) {conf: 0.30--0.50};
  \node[conf, right] at (6.5,-9.0) {conf: 0.20--0.40};

  % Entropy reduction
  \node[note, left] at (-6.7,0) {$H_0 \approx 8.0$ bit (tamamiyla bilinmiyor)};
  \node[note, left] at (-6.7,-1.8) {$\Delta H_1 \approx -3.5$ bit};
  \node[note, left] at (-6.7,-3.6) {$\Delta H_2 \approx -1.8$ bit};
  \node[note, left] at (-6.7,-5.4) {$\Delta H_3 \approx -1.0$ bit};
  \node[note, left] at (-6.7,-7.2) {$\Delta H_4 \approx -0.6$ bit};
  \node[note, left] at (-6.7,-9.0) {$\Delta H_5 \approx -0.3$ bit};

  % Arrows between layers
  \foreach \i/\j in {L1/L2, L2/L3, L3/L4, L4/L5, L5/L6} {
    \draw[arrow] (\i.south) -- (\j.north)
      node[midway, right, font=\tiny, text=gray!60!white] {isimsiz kalanlar};
  }

  % Title
  \node[above=0.8cm of L1, font=\large\bfseries, text=chapterblue]
    {Cascade Classifier --- 6 Katmanli Isimlendirme};
\end{tikzpicture}
\caption{CVariableNamer 6 katmanli cascade classifier mimarisi.
Her katman, bir onceki katmanin isimlendiremediyi sembolleri isler.
Sol tarafta bilgi entropisi ($H$) azalma miktarlari gosterilmektedir.}
\label{fig:c-namer-cascade}
\end{figure}

\subsubsection{Cascade Classifier Analojisi}

Bu mimari, Viola-Jones yuz tespit algoritmasindaki cascade classifier'a
benzer. Her katman (stage) bir ``filtre'' islevi gorur: belirli bir
guven seviyesinin ustundeki isimleri atar ve geri kalanlari bir sonraki
katmana iletir.

\begin{equation}
  P(\text{dogru isim}) = 1 - \prod_{k=1}^{6} \bigl(1 - p_k \cdot c_k\bigr)
  \label{eq:cascade-prob}
\end{equation}
burada $p_k$ katman $k$'nin yakalama olasiligi (recall) ve $c_k$ guven
skoru (confidence). Her katman bagimsiz bilgi kaynagi kullandigi icin,
toplam basari olasiligi kumulatif olarak artar.


\subsection{Katman 1: Symbol / Export Table}
\label{subsec:layer1-symbol}

En yuksek guven seviyesindeki ($c_1 = 0.95$) katmandir. Uc alt stratejisi vardir:

\begin{enumerate}
  \item \textbf{DWARF Debug Symbols:} Eger binary strip edilmemisse,
        \code{.debug\_info} section'indaki DIE (Debug Information Entry)
        yapilarindan orijinal fonksiyon, degisken ve tip isimleri dogrudan okunur.
  \item \textbf{Export Table (Mach-O/PE):} Dynamic linker icin tutulan export
        tablolari (\code{LC\_SYMTAB}, \code{\_\_LINKEDIT}) orijinal fonksiyon
        isimlerini icerir.
  \item \textbf{Objective-C RTTI:} \code{\_OBJC\_CLASS\_\$\_ClassName} pattern'i
        ile class isimleri, \code{sel\_refs} section'indan selector (method) isimleri
        elde edilir.
\end{enumerate}

\begin{lstlisting}[language=Python, caption={Ghidra otomatik isim tespiti regex pattern'leri (c\_namer.py)}, label={lst:ghidra-auto}]
_GHIDRA_AUTO_FUNC  = re.compile(r"^FUN_[0-9a-fA-F]+$")
_GHIDRA_AUTO_PARAM = re.compile(r"^param_\d+$")
_GHIDRA_AUTO_LOCAL = re.compile(r"^local_[0-9a-fA-F]+$")
_GHIDRA_AUTO_DAT   = re.compile(r"^DAT_[0-9a-fA-F]+$")
_GHIDRA_AUTO_PTR   = re.compile(r"^PTR_[0-9a-fA-F]+$")
_GHIDRA_AUTO_VAR   = re.compile(r"^[a-z]Var\d+$")   # uVar1, iVar2
_GHIDRA_AUTO_STACK = re.compile(r"^Stack\[0x[0-9a-fA-F]+\]$")
\end{lstlisting}

\begin{example}[Entropy Azaltma -- Katman 1]
Tipik bir macOS uygulamasinda 12.000 fonksiyon vardir. Bunlarin 3.000'i
(\%25) stripped degil ve dogrudan symbol table'dan alinir. Her fonksiyon
ismi icin bilgi entropisi:
\begin{equation}
  H_{\text{once}} = \log_2(N_{\text{olasi isimler}}) \approx \log_2(10^6) \approx 20 \text{ bit}
\end{equation}
Symbol table eslesmesinden sonra:
\begin{equation}
  H_{\text{sonra}} = 0 \text{ bit (kesin esleme)}
\end{equation}
Dolayisiyla $\Delta H_1 = -20$ bit, fakat bu sadece eslesen \%25 icin gecerlidir.
Agirliklai ortalama: $\Delta \bar{H}_1 = 0.25 \times (-20) = -5.0$ bit/fonksiyon.
\end{example}


\subsection{Katman 2: String Context Analysis}
\label{subsec:layer2-string}

Binary icindeki string literal'ler isim kurtarma icin zengin bir kaynaktir.
Log mesajlari, assert ifadeleri, hata metinleri ve UI string'leri genellikle
fonksiyon veya class isimlerini icerir.

\begin{tanim}[String Context Eslestirme]
Bir fonksiyon $f$ icindeki string referanslari $\{s_1, s_2, \ldots, s_m\}$
kumesi, $f$'in amacini gostermek icin analiz edilir:
\begin{itemize}
  \item \code{ClassName::MethodName} pattern'i $\Rightarrow$ $c = 0.85$
  \item \code{FuncName()} pattern'i $\Rightarrow$ $c = 0.80$
  \item Source path \code{/src/module/file.cpp} $\Rightarrow$ $c = 0.75$
  \item Steam/Valve keyword'leri $\Rightarrow$ $c = 0.65$
  \item Error verb extraction (``Failed to X'') $\Rightarrow$ $c = 0.55$
\end{itemize}
\end{tanim}

Keyword extraction sureci, string'den stopword'leri filtreleyerek anlamli
kelimeleri cikarir:

\begin{equation}
  \text{keywords}(s) = \text{filter}_{\neg\text{stopword}}\bigl(
    \text{tokenize}(s)\bigr) \cap \{w : |w| \geq 3\}
  \label{eq:keyword-extract}
\end{equation}

\file{c\_namer.py} icindeki \code{\_extract\_keywords\_from\_string} fonksiyonu
200+ stopword listesiyle filtreleme yapar ve sonucu en fazla 4 kelimeye sinirlar.


\subsection{Katman 3: API Call Pattern Matching}
\label{subsec:layer3-api}

Bir fonksiyonun cagirdigi sistem API'leri, fonksiyonun amacini guclu sekilde
gosterir. Ornegin \code{socket()}, \code{connect()}, \code{send()} uclusunu
cagiran bir fonksiyon, yuksek olasilikla bir ag istegi gonderme fonksiyonudur.

Kodda 60+ combo pattern tanimlanmistir:

\begin{longtable}{lll}
\toprule
\textbf{API Kombinasyonu} & \textbf{Atanan Isim} & \textbf{Conf.} \\
\midrule
\endhead
\code{\{socket, connect, send\}} & \code{send\_network\_request} & 0.80 \\
\code{\{fopen, fread, fclose\}} & \code{read\_file} & 0.80 \\
\code{\{socket, bind, listen, accept\}} & \code{start\_server\_socket} & 0.80 \\
\code{\{fork, exec\}} & \code{spawn\_process} & 0.80 \\
\code{\{malloc, memcpy, free\}} & \code{copy\_buffer} & 0.75 \\
\code{\{opendir, readdir, closedir\}} & \code{list\_directory} & 0.80 \\
\code{\{getaddrinfo\}} & \code{resolve\_address} & 0.70 \\
\code{\{strlen, strcpy\}} & \code{copy\_string} & 0.70 \\
\code{\{pthread\_create\}} & \code{create\_thread} & 0.70 \\
\code{\{mmap, munmap\}} & \code{map\_file\_memory} & 0.75 \\
\bottomrule
\caption{API combo pattern'lerinden bir secme (60+ pattern tanimli).}
\label{tab:api-combo}
\end{longtable}

Tek API cagrisi icin de pattern'ler vardir (200+ girdi):
\begin{lstlisting}[language=Python, caption={Tek API'den isim esleme ornekleri}, label={lst:single-api}]
_SINGLE_API_HINTS = {
    "CreateFileW":     "open_file_win",
    "RegOpenKeyExW":   "open_registry_key",
    "VirtualAlloc":    "allocate_virtual_memory",
    "NtQueryInformationProcess": "query_process_info",
    "dlopen":          "load_dynamic_library",
    ...
}
\end{lstlisting}


\subsection{Katman 4: Call Graph Structure}
\label{subsec:layer4-callgraph}

Fonksiyonun cagri grafindaki pozisyonu isim hakkinda bilgi verir.
Ghidra'nin \file{call\_graph.json} ciktisindaki caller/callee iliskileri
analiz edilir:

\begin{itemize}
  \item \textbf{Entry point fonksiyonlar:} Hicbir fonksiyon tarafindan
        cagirilmayan fonksiyonlar (\code{in-degree = 0}) $\Rightarrow$ \code{main}, \code{entry\_point}
  \item \textbf{Leaf fonksiyonlar:} Baska fonksiyon cagirmayan fonksiyonlar
        (\code{out-degree = 0}) $\Rightarrow$ \code{utility\_func}, \code{helper}
  \item \textbf{Hub fonksiyonlar:} Cok sayida farkli fonksiyondan cagirilir
        (\code{in-degree $\gg$ ortalama}) $\Rightarrow$ \code{common\_handler}, \code{dispatch}
  \item \textbf{Wrapper fonksiyonlar:} Tek bir fonksiyon cagirip donen
        (\code{out-degree = 1}, kisa govde) $\Rightarrow$ \code{wrapper\_XYZ}
\end{itemize}

\begin{equation}
  \text{role}(f) = \begin{cases}
    \text{entry} & \text{if } \deg^-(f) = 0 \\
    \text{leaf} & \text{if } \deg^+(f) = 0 \\
    \text{hub} & \text{if } \deg^-(f) > \mu^- + 2\sigma^- \\
    \text{wrapper} & \text{if } \deg^+(f) = 1 \wedge |f| < \tau
  \end{cases}
  \label{eq:callgraph-role}
\end{equation}
burada $\mu^-, \sigma^-$ in-degree ortalama ve standart sapma,
$|f|$ fonksiyon buyuklugu (instruction sayisi), $\tau$ esik degeridir.


\subsection{Katman 5: Dataflow Analysis}
\label{subsec:layer5-dataflow}

Parametre kullanimi, degisken atamalari ve veri akisi izlenerek isim cikarimi yapilir.
Bu katman ozellikle parametrelerin isimlendirilmesinde etkilidir.

\begin{itemize}
  \item \code{param\_1} bir \code{fopen()} cagrisinin ilk argumani olarak
        kullaniliyorsa $\Rightarrow$ \code{filename}
  \item \code{param\_2} bir \code{malloc()} cagrisinin argumani ise
        $\Rightarrow$ \code{size} veya \code{buffer\_size}
  \item \code{local\_28} bir dongu sayaciysa $\Rightarrow$ \code{i}, \code{loop\_counter}
\end{itemize}

Dataflow analizi, SSA (Static Single Assignment) benzeri bir yaklasimla
her degiskenin tum kullanim noktalarini (use-def chain) izler:

\begin{equation}
  \text{name}(v) = \text{argmax}_{n \in \mathcal{N}}
    \sum_{u \in \text{uses}(v)} w(u) \cdot \mathbb{1}[\text{context}(u) \Rightarrow n]
  \label{eq:dataflow-naming}
\end{equation}
burada $\mathcal{N}$ aday isimler kumesi, $\text{uses}(v)$ degisken $v$'nin
kullanildigi noktalar, $w(u)$ kullanim noktasinin bilgi agirligi ve
$\text{context}(u) \Rightarrow n$ kullanim baglaminin isim $n$'yi
destekleyip desteklemedigi.


\subsection{Katman 6: Type-Based Inference}
\label{subsec:layer6-type}

En dusuk guven seviyesindeki katmandir ($c_6 = 0.20$--$0.40$).
Ghidra'nin cikardigi tip bilgisinden isim cikarir:

\begin{itemize}
  \item \code{char *} tipi $\Rightarrow$ \code{str}, \code{name}, \code{path}
  \item \code{int} donus tipi + \code{bool} kullanim $\Rightarrow$ \code{is\_valid}, \code{check\_}
  \item \code{void *} parametresi $\Rightarrow$ \code{context}, \code{user\_data}
  \item Struct pointer $\Rightarrow$ struct adindan turetilmis isim
\end{itemize}

\subsection{Toplam Bilgi Entropisi Azaltma}
\label{subsec:total-entropy}

6 katmanin kumulatif etkisi, her sembol icin bilgi entropisi $H$'yi
sistematik olarak azaltir:

\begin{equation}
  H_{\text{final}} = H_0 - \sum_{k=1}^{6} \Delta H_k
  \label{eq:total-entropy}
\end{equation}

\begin{figure}[htbp]
\centering
\begin{tikzpicture}
\begin{axis}[
    width=12cm, height=7cm,
    ybar stacked,
    bar width=15pt,
    ylabel={Bilgi Entropisi (bit)},
    xlabel={Katman},
    symbolic x coords={Baslangic, K1, K2, K3, K4, K5, K6},
    xtick=data,
    ymin=0, ymax=22,
    legend style={at={(0.02,0.98)}, anchor=north west, fill=boxbg, text=white, draw=gray},
    axis background/.style={fill=black},
    every axis plot/.append style={fill opacity=0.85},
    axis line style={white},
    tick style={white},
    every tick label/.style={white},
    ylabel style={white},
    xlabel style={white},
]
\addplot[fill=red!60!black, draw=red!80] coordinates {
  (Baslangic,20) (K1,15) (K2,12.5) (K3,10.2) (K4,8.8) (K5,7.6) (K6,7.0)
};
\addlegendentry{Kalan Entropi}
\end{axis}
\end{tikzpicture}
\caption{6 katman boyunca bilgi entropisi azalma grafigi.
Baslangicta $H_0 \approx 20$ bit olan belirsizlik,
Katman 6 sonunda $H_6 \approx 7$ bit'e duser.}
\label{fig:entropy-reduction}
\end{figure}

\begin{example}[Sayisal Ornek: Steam Bootstrap Binary]
224 MB'lik bir macOS binary'sinde 48.000 fonksiyon tespit edilmistir:
\begin{itemize}
  \item Katman 1 (Symbol): 12.000 fonksiyon dogrudan isimlendirildi (\%25)
  \item Katman 2 (String): 8.400 ek fonksiyon (\%17.5)
  \item Katman 3 (API): 5.200 ek fonksiyon (\%10.8)
  \item Katman 4 (Call Graph): 3.600 ek fonksiyon (\%7.5)
  \item Katman 5 (Dataflow): 2.400 ek fonksiyon (\%5.0)
  \item Katman 6 (Type): 1.200 ek fonksiyon (\%2.5)
  \item \textbf{Toplam:} 32.800 / 48.000 = \textbf{\%68.3} isimlendirme basarisi
  \item Kalan 15.200 fonksiyon icin isim atanamaz (\code{FUN\_xxx} kalir)
\end{itemize}
\end{example}


% ============================================================================
\section{Tip Kurtarma (c\_type\_recoverer.py)}
\label{sec:type-recovery}

Ghidra'nin \code{undefined8}, \code{long *}, \code{void *} gibi generic
tipleri, \file{c\_type\_recoverer.py} tarafindan kullanim pattern'lerinden
analiz edilerek gercek struct, enum ve vtable tanimlarina donusturulur.

\subsection{Struct Recovery}

Ghidra decompile ciktisinda bir pointer'a offset'ler eklenerek farkli
field'lara erisim gosterilir:
\begin{lstlisting}[language=C, caption={Ghidra decompile ciktisinda struct field erisim pattern'i}]
*(long *)(param_1 + 0x10) = value;      // field @ offset 0x10
*(int *)(param_1 + 0x18) = 42;          // field @ offset 0x18
*(char **)(param_1 + 0x20) = "hello";   // field @ offset 0x20
\end{lstlisting}

Bu pattern'ler regex ile tespit edilir:

\begin{lstlisting}[language=Python, caption={Field access regex pattern'leri (c\_type\_recoverer.py)}, label={lst:field-access-regex}]
_FIELD_ACCESS_WRITE = re.compile(
    r"\*\(\s*"
    r"(?P<type>[a-zA-Z_][\w\s]*\*{0,2})\s*\*?\s*\)"
    r"\(\s*(?P<base>\w+)\s*\+\s*"
    r"(?P<offset>0x[0-9a-fA-F]+|\d+)\s*\)\s*="
)
\end{lstlisting}

Toplanan offset-tip ciftlerinden bir struct sentezlenir:

\begin{equation}
  \text{Struct}_i = \Bigl\{ (o_j, t_j, s_j) : j = 1, \ldots, n_i \Bigr\}
  \label{eq:struct-synth}
\end{equation}
burada $o_j$ byte offset, $t_j$ tip, $s_j$ byte boyutudur.
Alignment ARM64 icin 8 byte'dir.

\begin{example}[Struct Recovery Ornegi]
Asagidaki Ghidra ciktisi:
\begin{lstlisting}[language=C]
*(long *)(param_1 + 0x00) = uVar3;
*(int  *)(param_1 + 0x08) = 0;
*(char**)(param_1 + 0x10) = "default";
*(long *)(param_1 + 0x18) = 0x100;
\end{lstlisting}
Kurtarilan struct:
\begin{lstlisting}[language=C]
typedef struct {
    long   field_00;    // offset 0x00, 8 bytes
    int    field_08;    // offset 0x08, 4 bytes
    char  *field_10;    // offset 0x10, 8 bytes
    long   field_18;    // offset 0x18, 8 bytes
} recovered_struct_001; // total: 32 bytes, alignment: 8
\end{lstlisting}
\end{example}

\subsection{Enum Recovery}

Switch-case yapilari ve karsilastirma pattern'leri enum tespiti icin kullanilir:

\begin{lstlisting}[language=Python, caption={Switch/case regex pattern'leri}]
_SWITCH_START = re.compile(r"switch\s*\(\s*(?P<var>\w+)\s*\)")
_CASE_VALUE   = re.compile(r"case\s+(?P<val>-?\d+|0x[0-9a-fA-F]+)\s*:")
\end{lstlisting}

Bir switch ifadesinde 4'ten fazla case degeri varsa, bu degerler bir
enum olarak sentezlenir:

\begin{equation}
  \text{Enum}_i = \Bigl\{ (\text{name}_j, v_j) : j = 1, \ldots, m_i \Bigr\}
  \label{eq:enum-synth}
\end{equation}

\subsection{VTable Recovery}

Fonksiyon pointer atamalari vtable yapisini gosterir:
\begin{lstlisting}[language=C]
*(code **)(obj + 0x00) = FUN_001234;   // virtual method 0
*(code **)(obj + 0x08) = FUN_005678;   // virtual method 1
*(code **)(obj + 0x10) = FUN_009abc;   // virtual method 2
\end{lstlisting}

Regex: \code{\_VTABLE\_ASSIGN} pattern'i \code{code **} tipli pointer
atamalari\-ni yakalar ve \code{(offset, fonksiyon\_adi, signature)} uclulerini toplar.


% ============================================================================
\section{Algoritma Tespiti (c\_algorithm\_id.py)}
\label{sec:algorithm-id}

Decompile edilmis C kodunda kullanilan kriptografik algoritmalari, hash
fonksiyonlarini ve diger bilinen algoritmalari tespit eder.
Uc katmanli bir yaklasim kullanilir.

\subsection{Katman 1: Sabit-Bazli Tanimlar (Constant-Based)}

Kriptografik algoritmalar, kolayca degistirilemeyen matematiksel sabitlere
sahiptir. Bu sabitlerin kodda varligi, algoritmanin kullanildiginin guclu
kanit\-idir.

\begin{longtable}{llp{6cm}c}
\toprule
\textbf{Algoritma} & \textbf{Sabit Adi} & \textbf{Ilk Byte'lar (hex)} & \textbf{Conf.} \\
\midrule
\endhead
AES & S-box & \code{63 7C 77 7B F2 6B 6F C5} & 0.95 \\
AES & Inverse S-box & \code{52 09 6A D5 30 36 A5 38} & 0.95 \\
AES & Rcon & \code{01 02 04 08 10 20 40 80} & 0.90 \\
SHA-256 & K constants & \code{428A2F98 71374491 B5C0FBCF E9B5DBA5} & 0.95 \\
SHA-256 & H init & \code{6A09E667 BB67AE85 3C6EF372 A54FF53A} & 0.95 \\
SHA-512 & K constants & \code{428A2F98D728AE22 7137449123EF65CD} & 0.95 \\
SHA-1 & H init & \code{67452301 EFCDAB89 98BADCFE 10325476} & 0.95 \\
MD5 & T values & \code{D76AA478 E8C7B756 242070DB C1BDCEEE} & 0.95 \\
CRC32 & Table start & \code{00000000 77073096 EE0E612C 990951BA} & 0.90 \\
ChaCha20 & Sigma & \code{61707865 3320646E 79622D32 6B206574} & 0.95 \\
Blowfish & P-array & \code{243F6A88 85A308D3 13198A2E 03707344} & 0.90 \\
HMAC & ipad/opad & \code{36 36 36 36} / \code{5C 5C 5C 5C} & 0.80 \\
Poly1305 & Clamp & \code{0FFFFFFF 0FFFFFFC 0FFFFFFC 0FFFFFFC} & 0.90 \\
\bottomrule
\caption{Sabit-bazli algoritma tespiti icin kullanilan kriptografik sabitlerin bir bolumu.
Kodda toplam 25+ algoritma icin 40+ sabit grubu tanimlidir.}
\label{tab:crypto-constants}
\end{longtable}

\begin{tanim}[AES S-Box Tespiti]
AES S-box, Galois Field $GF(2^8)$'de carpimsal ters alma ve affine
donusum ile hesaplanan 256 byte'lik bir lookup table'dir:
\begin{equation}
  S(x) = A \cdot x^{-1} \oplus c, \quad x \in GF(2^8), \quad
  \text{irreducible polynomial: } m(x) = x^8 + x^4 + x^3 + x + 1
  \label{eq:aes-sbox}
\end{equation}
Ilk 8 byte'in kodda bulunmasi yeterli kanittir ($p_{\text{false positive}} < 10^{-15}$).
\end{tanim}

\subsection{Katman 2: Yapi-Bazli Tespit (Structure-Based)}

Kod yapisindan algoritma pattern'leri taninir. 14 yapisal pattern tanimlidir:

\begin{itemize}
  \item \textbf{Feistel Network:} $L_{i+1} = R_i$, $R_{i+1} = L_i \oplus F(R_i, K_i)$
        pattern'i -- DES, Blowfish icin.
  \item \textbf{S-box Lookup:} \code{table[x \& 0xFF]} veya \code{table[(x >> n) \& 0xFF]}
        -- AES SubBytes icin.
  \item \textbf{Hash Rounds:} Rotate (\code{>>>}, \code{<<<}),
        XOR ve toplama islemlerinin 16+ kez tekrari -- SHA, MD5 icin.
  \item \textbf{Galois Field:} Moduler aritmetik (\code{a * b \% p}) -- ECC icin.
  \item \textbf{Key Schedule:} \code{round\_key}, \code{key\_schedule} degiskenleri.
  \item \textbf{Base64 Encoding:} Alfabe string'i \code{A-Za-z0-9+/}.
\end{itemize}

\subsection{Katman 3: API Korelasyonu}

Bilinen crypto API cagrilari dogrudan algoritma tespiti saglar.
Kodda 50+ crypto API tanimlidir:

\begin{itemize}
  \item \textbf{Apple CommonCrypto:} \code{CCCrypt}, \code{CC\_SHA256}, \code{CCHmac}
  \item \textbf{OpenSSL:} \code{EVP\_EncryptInit}, \code{SSL\_read}, \code{RSA\_public\_encrypt}
  \item \textbf{Windows CryptoAPI:} \code{CryptEncrypt}, \code{BCryptOpenAlgorithmProvider}
  \item \textbf{Sodium:} \code{crypto\_secretbox}, \code{crypto\_aead\_}
\end{itemize}


% ============================================================================
\section{Yorum Uretimi (c\_comment\_generator.py)}
\label{sec:comment-gen}

Decompile edilmis C dosyalarina bes tip yorum eklenir:

\begin{enumerate}
  \item \textbf{Function Header Block:} Doxygen benzeri fonksiyon baslik yorumu.
        Parametre tipleri, donus degeri, cagiran/cagrilan fonksiyonlar, tespit
        edilen algoritmalar.
  \item \textbf{System Call Annotations:} 100+ sistem cagrisi icin aciklama.
        Ornegin: \code{malloc(param\_1)} $\Rightarrow$ \code{/* allocates \{arg0\} bytes on heap */}
  \item \textbf{Vulnerability Annotations:} Tehlikeli fonksiyon kullanimlari icin
        uyari. \code{strcpy} $\Rightarrow$ \code{/* WARNING: no bounds check, use strncpy */}
  \item \textbf{Control Flow Annotations:} State machine, infinite loop,
        buyuk switch-case yapilari icin aciklama.
  \item \textbf{Algorithm Labels:} \code{c\_algorithm\_id} sonuclarindan
        algoritma etiketleri.
\end{enumerate}

\begin{lstlisting}[language=C, caption={Otomatik yorum uretimi ornegi}]
/**
 * Function: send_network_request (FUN_00102a40)
 * @brief Sends data over network socket
 * @param param_1 (renamed: socket_fd) - socket file descriptor
 * @param param_2 (renamed: buffer) - data buffer pointer
 * @param param_3 (renamed: length) - buffer length
 * @return int - bytes sent or -1 on error
 *
 * Calls: socket(), connect(), send(), close()
 * Called by: FUN_00103b80, FUN_00104c20
 * Algorithm: None detected
 * Confidence: 0.80 (API pattern)
 */
int send_network_request(int socket_fd, void *buffer, long length) {
    /* socket(AF_INET, SOCK_STREAM, 0) - creates network socket */
    int fd = socket(2, 1, 0);
    ...
}
\end{lstlisting}


% ============================================================================
\section{Proje Yapilandirma}
\label{sec:project-structure}

\subsection{project\_reconstructor.py -- Modul Iskele Olusturma}

JavaScript bundle'lardan elde edilen modullerin gercek bir proje yapisina
donusturulmesini saglar. 8 adimlik pipeline:

\begin{enumerate}
  \item Modul dizinini bul (\code{deep\_modules/} veya \code{webpack\_modules/})
  \item \code{reconstruct-project.mjs} script'i ile 30+ kategori pattern'ine
        gore dosya yerlestirme
  \item \code{DependencyResolver} ile gercek npm paketlerini tespit
  \item \code{package.json} olusturma
  \item Entry point belirleme
  \item README olusturma
  \item \code{.gitignore} olusturma
  \item Proje yapisini dogrulama (validation)
\end{enumerate}

\subsection{C Project Builder (c\_project\_builder.py)}

Decompile edilmis binary'den organize bir IDE projesi uretir. Subsystem
siniflandirma kurallari ile fonksiyonlar mantiksal gruplara ayrilir:

\begin{figure}[htbp]
\centering
\begin{tikzpicture}[
    box/.style={draw=sectioncyan, fill=boxbg, text=white, minimum width=3cm,
                minimum height=0.8cm, rounded corners=2pt, font=\small},
    dir/.style={draw=yellow!70!white, fill=boxbg, text=yellow!70!white,
                minimum width=2.5cm, minimum height=0.7cm, font=\small\ttfamily},
    conn/.style={-{Stealth[length=4pt]}, white, thick},
]
  % Root
  \node[box] (root) at (0,0) {project/};

  % Level 1
  \node[dir] (src) at (-4,-1.5) {src/};
  \node[dir] (inc) at (0,-1.5) {include/};
  \node[dir] (ana) at (4,-1.5) {analysis/};

  % src children
  \node[dir] (main) at (-6,-3) {main.c};
  \node[dir] (net) at (-4,-3) {network/};
  \node[dir] (cry) at (-2,-3) {crypto/};
  \node[dir] (fio) at (-4,-4.5) {file\_io/};
  \node[dir] (misc) at (-2,-4.5) {misc/};

  % include children
  \node[dir] (types) at (0,-3) {types.h};

  % analysis children
  \node[dir] (cfg) at (3,-3) {call\_graph.dot};
  \node[dir] (algo) at (5,-3) {algorithms.md};

  % Files
  \node[dir] (cmake) at (0,-5.5) {CMakeLists.txt};

  % Connections
  \draw[conn] (root) -- (src);
  \draw[conn] (root) -- (inc);
  \draw[conn] (root) -- (ana);
  \draw[conn] (src) -- (main);
  \draw[conn] (src) -- (net);
  \draw[conn] (src) -- (cry);
  \draw[conn] (src) -- (fio);
  \draw[conn] (src) -- (misc);
  \draw[conn] (inc) -- (types);
  \draw[conn] (ana) -- (cfg);
  \draw[conn] (ana) -- (algo);
  \draw[conn] (root) -- (cmake);
\end{tikzpicture}
\caption{C Project Builder cikti yapisi. Fonksiyonlar API kullanim pattern'lerine
gore subsystem klasorlerine (network, crypto, file\_io) dagitilir.}
\label{fig:c-project-structure}
\end{figure}

\subsection{project\_builder.py -- CMake/Makefile Uretimi}

JavaScript projeler icin \code{npm install \&\& npm start} ile calisabilecek
bir proje uretir:

\begin{itemize}
  \item \code{package.json}: Otomatik tespit edilen dependency'ler, script'ler
  \item \code{.env.example}: Ortam degiskenleri sablonu
  \item \code{Dockerfile}: Multi-stage build ile production-ready container
  \item \code{README.md}: Proje dokumantasyonu
\end{itemize}


% ============================================================================
\section{NPM Source Matching}
\label{sec:npm-source-matching}

Minified JavaScript bundle'larindaki fonksiyonlari orijinal npm paketi
kaynak koduyla eslestirme.

\subsection{AST Structural Hash (ast\_fingerprinter.py)}

Her fonksiyon icin yapisal parmak izi cikarilir. Minifier'lar isimleri
degistirir ama yapisal ozellikleri korur:

\begin{equation}
  \text{FP}(f) = \bigl(\text{arity}(f),\; \mathcal{S}(f),\; \mathcal{P}(f),\;
    B(f),\; L(f),\; R(f)\bigr)
  \label{eq:function-fingerprint}
\end{equation}
burada:
\begin{itemize}
  \item $\text{arity}(f)$: Parametre sayisi
  \item $\mathcal{S}(f)$: String literal'ler kumesi (minifier bunlara dokunmaz)
  \item $\mathcal{P}(f)$: Property access'ler kumesi (\code{.push}, \code{.map})
  \item $B(f)$: Dallanma yapisi (if/switch/ternary sayilari)
  \item $L(f)$: Dongu yapisi (for/while/do-while sayilari)
  \item $R(f)$: Return/throw sayilari
\end{itemize}

\subsection{Tree Edit Distance (Zhang-Shasha)}

Iki fonksiyonun yapisal benzerligi, AST'lerin tree edit distance'i
ile hesaplanir:

\begin{equation}
  d_{\text{edit}}(T_1, T_2) = \min_{\pi \in \Pi} \sum_{i} c(\pi_i)
  \label{eq:tree-edit-distance}
\end{equation}
burada $\Pi$ tum gecerli edit islem dizileri kumesi,
$c(\pi_i)$ her islemin maliyeti (insert, delete, relabel).

Zhang-Shasha algoritmasi $\bigO{|T_1| \cdot |T_2| \cdot \min(d_1, l_1) \cdot \min(d_2, l_2)}$
zaman karmasikligina sahiptir; burada $d_i$ derinlik ve $l_i$ yaprak sayisi.

Normallestirilmis benzerlik:
\begin{equation}
  \text{sim}(T_1, T_2) = 1 - \frac{d_{\text{edit}}(T_1, T_2)}{|T_1| + |T_2|}
  \label{eq:normalized-similarity}
\end{equation}

\subsection{Structural Matcher Pipeline}

\begin{figure}[htbp]
\centering
\begin{tikzpicture}[
    pstep/.style={draw=sectioncyan, fill=boxbg, text=white,
                 minimum width=4cm, minimum height=1cm,
                 rounded corners=3pt, font=\small},
    arrow/.style={-{Stealth[length=5pt]}, thick, white},
]
  \node[pstep] (s1) at (0,0) {1. Fingerprint Extraction};
  \node[pstep] (s2) at (0,-1.8) {2. Pairwise Similarity};
  \node[pstep] (s3) at (0,-3.6) {3. Greedy Matching};
  \node[pstep] (s4) at (0,-5.4) {4. Consistency Check};
  \node[pstep] (s5) at (0,-7.2) {5. Name Transfer};

  \draw[arrow] (s1) -- (s2);
  \draw[arrow] (s2) -- (s3);
  \draw[arrow] (s3) -- (s4);
  \draw[arrow] (s4) -- (s5);

  \node[right=1cm, text=gray!70!white, font=\scriptsize, align=left] at (s1.east)
    {ASTFingerprinter: regex-based\\fonksiyon sinir tespiti};
  \node[right=1cm, text=gray!70!white, font=\scriptsize, align=left] at (s2.east)
    {$O(n \times m)$ cift karsilastirma\\similarity matrix olustur};
  \node[right=1cm, text=gray!70!white, font=\scriptsize, align=left] at (s3.east)
    {En yuksek benzerlikten basla\\eslesen ciftleri cikar};
  \node[right=1cm, text=gray!70!white, font=\scriptsize, align=left] at (s4.east)
    {Toplu tutarlilik kontrolu\\coverage $\geq$ consistency\_ratio};
  \node[right=1cm, text=gray!70!white, font=\scriptsize, align=left] at (s5.east)
    {Orijinal isimler minified koda\\transfer edilir};
\end{tikzpicture}
\caption{Source matching pipeline -- 5 adimda minified fonksiyonlarin
orijinal isimleri kurtarilir.}
\label{fig:source-matching-pipeline}
\end{figure}

Greedy matching algoritmasi:
\begin{enumerate}
  \item Her (minified, orijinal) cifti icin $\text{sim}(f_i^{\text{min}}, f_j^{\text{orig}})$ hesapla
  \item Tum ciftleri azalan benzerlige gore sirala
  \item En yuksekten baslayarak eslestir (matched olanlari cikar)
  \item $\text{sim} < \text{min\_similarity}$ (varsayilan $0.65$) alti olanlarini at
\end{enumerate}


% ============================================================================
\section{Parametre Kurtarma (param\_recovery.py)}
\label{sec:param-recovery}

Obfuscated JavaScript fonksiyonlarindaki tek harfli parametre isimlerini
anlamli isimlere donusturur. 5 strateji kullanilir:

\begin{longtable}{llcc}
\toprule
\textbf{Strateji} & \textbf{Pattern} & \textbf{Confidence} & \textbf{Oncelik} \\
\midrule
\endhead
this.X = param & \code{this.name = a} $\Rightarrow$ \code{a} $\to$ \code{name} & 0.90--0.95 & 1 \\
Destructuring & \code{const \{X,Y\} = a} $\Rightarrow$ \code{a} $\to$ \code{options} & 0.88 & 2 \\
Call-site object & \code{fn(\{key: val\})} $\Rightarrow$ param $\to$ \code{options} & 0.82 & 3 \\
Call-site identifier & \code{fn(userId)} $\Rightarrow$ param $\to$ \code{userId} & 0.72 & 4 \\
Property access & \code{a.length} $\Rightarrow$ \code{a} $\to$ \code{array} & 0.35--0.78 & 5 \\
\bottomrule
\caption{Parametre kurtarma stratejileri ve guven seviyeleri.}
\label{tab:param-recovery}
\end{longtable}

\subsection{TypeScript .d.ts Mapping (dts\_namer.py)}

TypeScript tanim dosyalarindan export isimleri cikarilir ve minified
koddaki export'larla eslestirilir:

\begin{enumerate}
  \item \textbf{Kaynak stratejisi:}
  \begin{itemize}
    \item DefinitelyTyped: \code{@types/xxx} paketi (npm registry)
    \item Paketin kendi \code{.d.ts} dosyasi (\code{package.json} ``types'' alani)
    \item \code{unpkg.com} CDN uzerinden indirme (fallback)
  \end{itemize}
  \item \textbf{Eslestirme stratejisi:} Export sirasi build tool'lar tarafindan
        korunur; isim uzunluklari ve tip bilgileri ile cross-validation.
\end{enumerate}

\begin{equation}
  \text{match}(e_i^{\text{min}}, e_j^{\text{dts}}) =
  \begin{cases}
    1 & \text{if } \text{kind}_i = \text{kind}_j \wedge |\text{arity}_i - \text{arity}_j| \leq 1 \\
    0 & \text{otherwise}
  \end{cases}
  \label{eq:dts-match}
\end{equation}


\begin{ozet}
\begin{itemize}
  \item Kaynak kod rekonstruksiyonu, 6 katmanli cascade classifier ile fonksiyon
        isimlendirme yapar (symbol $\to$ string $\to$ API $\to$ call graph
        $\to$ dataflow $\to$ type).
  \item Her katman bagimsiz bilgi kaynagi kullanir ve toplam basari kumulatif artar.
  \item Tip kurtarma, field access pattern'lerinden struct, switch/case'den enum,
        function pointer'lardan vtable sentezler.
  \item Algoritma tespiti 3 katmanli: sabit-bazli (AES S-box, SHA-256 K),
        yapi-bazli (Feistel, hash rounds) ve API korelasyonu.
  \item NPM source matching, AST fingerprint ve greedy matching ile minified
        fonksiyonlarin orijinal isimlerini kurtarir.
  \item Parametre kurtarma 5 strateji kullanir: this.X=param en yuksek guvenle.
\end{itemize}
\end{ozet}


\newpage

% ############################################################################
%
%                    BOLUM 8 -- GUVENLIK ANALIZ ALGORITMALARI
%
% ############################################################################

\section{Guvenlik Analiz Algoritmalarinin Temeli}
\label{sec:security-giris}

Karadul'un guvenlik analiz katmani, binary ve bytecode duzeyinde
tehditleri tespit etmek ve koruma mekanizmalarini asmak icin
bir dizi algoritmik teknik kullanir. Bu bolumde bu algoritmalarin
matematiksel temelleri, karmasikliklari ve implementasyon detaylari
ele alinir.


% ============================================================================
\section{Shannon Entropy ve Packer Tespiti}
\label{sec:entropy-packer}

\subsection{Shannon Entropy Formulu}

Bir rasgele degisken $X$'in bilgi entropisi (Shannon, 1948):

\begin{equation}
  H(X) = -\sum_{x \in \mathcal{X}} p(x) \log_2 p(x)
  \label{eq:shannon-entropy}
\end{equation}

\begin{tanim}[Entropi Turetmesi]
Bilgi entropisi, bir mesajin ortalama bilgi icerigi olarak turetilebilir.
$N$ byte'lik bir binary dosyada $x_i$ byte degerinin frekans\-i $f_i$ ise,
olasilik tahmini:
\begin{equation}
  \hat{p}(x_i) = \frac{f_i}{N}, \quad i = 0, 1, \ldots, 255
  \label{eq:byte-prob}
\end{equation}
Binary verilerin entropisi $[0, 8]$ bit araligindadir (byte basina):
\begin{itemize}
  \item $H = 0$: Tum byte'lar ayni (ornegin tamamen sifir)
  \item $H = 8$: Tum byte'lar esit olasilikli (uniform rastgele, sifreli veya sikistirilmis)
\end{itemize}
\end{tanim}

\begin{lstlisting}[language=Python, caption={Shannon entropy implementasyonu (packed\_binary.py)}, label={lst:entropy-impl}]
def calculate_entropy(data: bytes) -> float:
    if not data:
        return 0.0
    freq = [0] * 256
    for byte in data:
        freq[byte] += 1
    length = len(data)
    entropy = 0.0
    for f in freq:
        if f > 0:
            p = f / length
            entropy -= p * math.log2(p)
    return entropy
\end{lstlisting}

Algoritma karmasikligi: $\bigO{N + 256} = \bigO{N}$ -- dosya boyutu $N$ ile dogrusal.

\subsection{Section-Level Entropy Profiling}

Tum dosyanin entropisi yerine, section'lara bolup her birinin entropisini
hesaplamak daha bilgilendiricidir:

\begin{equation}
  H_j = H(X_j), \quad j = 1, \ldots, \lceil N / S \rceil
  \label{eq:section-entropy}
\end{equation}
burada $S$ section boyutu (varsayilan 65536 byte = 64 KB).

\begin{lstlisting}[language=Python, caption={Section-level entropy (packed\_binary.py)}]
def calculate_section_entropy(data: bytes, section_size: int = 65536):
    entropies = []
    for i in range(0, len(data), section_size):
        chunk = data[i:i + section_size]
        if len(chunk) >= 256:  # cok kucuk chunk yaniltici
            entropies.append(calculate_entropy(chunk))
    return entropies
\end{lstlisting}

\subsection{Packing Tespit Karari}

\file{PackingDetector} sinifi, 5 adimlik bir tespit pipeline'i kullanir:

\begin{figure}[htbp]
\centering
\begin{tikzpicture}[
    decision/.style={diamond, draw=warningborder, fill=warningbg, text=white,
                     minimum width=2cm, minimum height=1.2cm, font=\small,
                     inner sep=2pt, aspect=2},
    process/.style={draw=sectioncyan, fill=boxbg, text=white,
                    minimum width=3.5cm, minimum height=0.8cm,
                    rounded corners=2pt, font=\small},
    result/.style={draw=successborder, fill=successbg, text=white,
                   minimum width=3cm, minimum height=0.7cm,
                   rounded corners=2pt, font=\small\bfseries},
    arrow/.style={-{Stealth[length=4pt]}, thick, white},
    yes/.style={font=\scriptsize, text=green!70!white},
    no/.style={font=\scriptsize, text=red!70!white},
]
  \node[process] (start) at (0,0) {Binary Oku};
  \node[decision] (upx) at (0,-1.8) {UPX magic?};
  \node[result] (r1) at (4,-1.8) {UPX (c=0.95)};

  \node[decision] (py) at (0,-3.8) {PyInstaller?};
  \node[result] (r2) at (4,-3.8) {PyInstaller (c=0.95)};

  \node[decision] (nu) at (0,-5.8) {Nuitka?};
  \node[result] (r3) at (4,-5.8) {Nuitka (c=0.85)};

  \node[decision] (ent) at (0,-8) {$H > 7.0$ ?};
  \node[decision] (imp) at (0,-10.2) {imports $< 10$ ?};
  \node[result] (r4) at (4,-8) {Unknown Packed};
  \node[result] (r5) at (4,-10.2) {Generic Packed};
  \node[result] (r6) at (0,-12) {Not Packed};

  \draw[arrow] (start) -- (upx);
  \draw[arrow] (upx) -- (r1) node[midway, yes, above] {evet};
  \draw[arrow] (upx) -- (py) node[midway, no, left] {hayir};
  \draw[arrow] (py) -- (r2) node[midway, yes, above] {evet};
  \draw[arrow] (py) -- (nu) node[midway, no, left] {hayir};
  \draw[arrow] (nu) -- (r3) node[midway, yes, above] {evet};
  \draw[arrow] (nu) -- (ent) node[midway, no, left] {hayir};
  \draw[arrow] (ent) -- (r4) node[midway, yes, above] {evet + ratio $>$ 0.5};
  \draw[arrow] (ent) -- (imp) node[midway, no, left] {hayir};
  \draw[arrow] (imp) -- (r5) node[midway, yes, above] {$H > 6.5$};
  \draw[arrow] (imp) -- (r6) node[midway, no, left] {hayir};
\end{tikzpicture}
\caption{PackingDetector karar agaci. Her adimda farkli bir kanit kaynagi
kontrol edilir.}
\label{fig:packing-decision-tree}
\end{figure}

\subsection{Packer Signature'lari}

\begin{longtable}{llcc}
\toprule
\textbf{Packer} & \textbf{Tespit Yontemi} & \textbf{Magic Bytes} & \textbf{Conf.} \\
\midrule
\endhead
UPX & Magic byte \code{UPX!} & \code{55 50 58 21} & 0.95 \\
PyInstaller & MEI magic (8 byte) & \code{4D 45 49 0C 0B 0A 0B 0E} & 0.95 \\
Nuitka & String patterns & \code{\_\_nuitka\_} + 4 variant & 0.85 \\
Themida & High entropy + VM stub & Entropy $> 7.5$ & 0.75 \\
VMProtect & Virtualized sections & Section name + entropy & 0.70 \\
\bottomrule
\caption{Desteklenen packer tipleri ve tespit yontemleri.}
\label{tab:packer-sigs}
\end{longtable}

\subsection{Chi-Squared Uniformity Test}

Section entropisine ek olarak, byte dagiliminin uniform olup olmadigini
test etmek icin chi-squared testi kullanilir:

\begin{equation}
  \chi^2 = \sum_{i=0}^{255} \frac{(O_i - E_i)^2}{E_i}, \quad
  E_i = \frac{N}{256}
  \label{eq:chi-squared}
\end{equation}
burada $O_i$ gozlenen frekans, $E_i$ beklenen frekans.
$\chi^2 < 293.25$ (\%99.9 guven arali\-gi, $df=255$) ise byte dagilimi
uniform kabul edilir -- bu da guclu sikilasma veya sifreleme isaretidir.

\begin{example}[Sayisal Ornek]
1 MB'lik bir section'da:
\begin{itemize}
  \item Her byte degeri $E = 1048576/256 = 4096$ kez beklenir
  \item Uniform rastgele veride: $\chi^2 \approx 255$ (serbestlik derecesi)
  \item Normal C kodunda: $\chi^2 \approx 50000$+ (0x00 ve 0xFF dominant)
  \item Packed veride: $\chi^2 \approx 200$--$300$ (uniform'a yakin)
\end{itemize}
\end{example}


% ============================================================================
\section{YARA Pattern Matching}
\label{sec:yara}

\subsection{Aho-Corasick Algoritma Temeli}

YARA, pattern eslestirme icin Aho-Corasick coklu pattern arama
algoritmasini kullanir:

\begin{equation}
  T(n, m, z) = \bigO{n + m + z}
  \label{eq:aho-corasick}
\end{equation}
burada $n$ metin uzunlugu, $m$ tum pattern'lerin toplam uzunlugu,
$z$ esleme sayisi.

\begin{tanim}[Aho-Corasick Otomatonu]
Aho-Corasick, pattern kumesi $P = \{p_1, p_2, \ldots, p_k\}$ icin
bir trie (prefix tree) olusturur ve bu trie'yi failure link'lerle
zenginlestirir:
\begin{enumerate}
  \item \textbf{Goto fonksiyonu} $g(s, a)$: State $s$'den karakter $a$
        ile sonraki state.
  \item \textbf{Failure fonksiyonu} $f(s)$: Esleme basarisiz oldugunda
        geri donus state'i (KMP'nin failure fonksiyonunun genellestirilmesi).
  \item \textbf{Output fonksiyonu} $\text{out}(s)$: State $s$'de tamamlanan
        pattern'ler.
\end{enumerate}
\end{tanim}

\begin{figure}[htbp]
\centering
\begin{tikzpicture}[
    state/.style={circle, draw=sectioncyan, fill=boxbg, text=white,
                  minimum size=0.8cm, font=\small},
    accept/.style={state, double, double distance=2pt},
    trans/.style={-{Stealth[length=4pt]}, thick, white},
    fail/.style={-{Stealth[length=4pt]}, dashed, red!60!white, bend left=30},
    lbl/.style={font=\scriptsize, text=cyan},
]
  \node[state] (0) at (0,0) {0};
  \node[state] (1) at (2,1) {1};
  \node[state] (2) at (4,1) {2};
  \node[accept] (3) at (6,1) {3};
  \node[state] (4) at (2,-1) {4};
  \node[accept] (5) at (4,-1) {5};

  \draw[trans] (0) -- (1) node[midway, lbl, above] {\code{U}};
  \draw[trans] (1) -- (2) node[midway, lbl, above] {\code{P}};
  \draw[trans] (2) -- (3) node[midway, lbl, above] {\code{X}};
  \draw[trans] (0) -- (4) node[midway, lbl, below] {\code{M}};
  \draw[trans] (4) -- (5) node[midway, lbl, below] {\code{E}};

  \draw[fail] (2) to (0);
  \draw[fail] (5) to (0);

  \node[right=0.5cm of 3, text=green!70!white, font=\scriptsize] {output: ``UPX''};
  \node[right=0.5cm of 5, text=green!70!white, font=\scriptsize] {output: ``ME''};
\end{tikzpicture}
\caption{Basitlestirilmis Aho-Corasick otomatonu ornegi: ``UPX'' ve ``ME''
pattern'leri icin. Kesikli oklar failure link'lerdir.}
\label{fig:aho-corasick}
\end{figure}

\subsection{YARA Kural Yapisi}

Karadul, YARA kurulu degilse regex-based fallback ile calisir.
40+ dahili kural tanimlanmistir:

\begin{lstlisting}[language=Python, caption={Dahili YARA kural ornekleri (yara\_scanner.py)}]
rules.append(BuiltinRule(
    name="AES_SBox",
    tags=["crypto", "aes"],
    meta={"description": "AES S-Box lookup table"},
    byte_patterns=[
        _hex("63 7C 77 7B F2 6B 6F C5 30 01 67 2B FE D7 AB 76"),
    ],
))

rules.append(BuiltinRule(
    name="SHA256_Init_Constants",
    tags=["crypto", "sha256"],
    byte_patterns=[
        _hex("6A09E667 BB67AE85 3C6EF372 A54FF53A"),
    ],
))
\end{lstlisting}

\subsection{Hex Pattern + Wildcard}

YARA hex pattern'leri wildcard (\code{??}) ve jump (\code{[n-m]}) destekler:
\begin{equation}
  \text{pattern}: \quad \code{63 7C ?? 7B [2-4] 6B 6F C5}
  \label{eq:yara-wildcard}
\end{equation}
Bu, ``\code{63 7C}'' ile basla, herhangi bir byte, \code{7B}, 2 ila 4
byte atlama, sonra ``\code{6B 6F C5}'' seklinde eslesir. Bu esneklik
compiler optimization farklarini absorbe eder.


% ============================================================================
\section{Control Flow Graph (CFG) Analizi}
\label{sec:cfg-analysis}

\subsection{Formal Tanim}

\begin{tanim}[Control Flow Graph]
Bir programin CFG'si $G = (V, E, v_{\text{entry}}, v_{\text{exit}})$
yonlu grafi ile tanimlanir:
\begin{itemize}
  \item $V$: Basic block'lar kumesi
  \item $E \subseteq V \times V$: Kontrol akisi kenarlari
  \item $v_{\text{entry}} \in V$: Giris blogu
  \item $v_{\text{exit}} \in V$: Cikis blogu
\end{itemize}
\end{tanim}

\subsection{Basic Block ve Leader Detection}

Bir \emph{basic block}, kontrol akisinin yalnizca basindan girdigi ve
yalnizca sonundan ciktigi maksimal instruction dizisidir.

\emph{Leader} tespiti kurallari:
\begin{enumerate}
  \item Ilk instruction bir leader'dir
  \item Bir dallanma (jump) hedefi olan instruction bir leader'dir
  \item Bir dallanma instruction'indan hemen sonraki instruction bir leader'dir
\end{enumerate}

\subsection{Dominator Tree -- Lengauer-Tarjan}

\begin{tanim}[Dominator]
$d \in V$ dugumleri $n \in V$ dugumunu \emph{dominate} eder ($d \text{ dom } n$)
eger ve ancak $v_{\text{entry}}$'den $n$'ye giden her yol $d$'den geciyorsa.
\end{tanim}

\begin{equation}
  \text{Dom}(n) = \{n\} \cup \bigcap_{p \in \text{pred}(n)} \text{Dom}(p)
  \label{eq:dominator-set}
\end{equation}

Lengauer-Tarjan algoritmasi dominator tree'yi neredeyse dogrusal zamanda hesaplar:

\begin{equation}
  T_{\text{LT}}(n) = \bigO{n \cdot \alpha(n)}
  \label{eq:lengauer-tarjan}
\end{equation}
burada $\alpha(n)$ ters Ackermann fonksiyonudur (pratikte $\leq 4$, yani
neredeyse $\bigO{n}$).

\subsection{Loop Detection ve Back Edges}

\begin{tanim}[Back Edge]
Bir kenar $(u, v) \in E$, eger $v$ dugumu $u$'yu dominate ediyorsa
bir \emph{back edge}'dir. Her back edge bir natural loop tanimlar.
\end{tanim}

Natural loop tespit algoritmasi:
\begin{enumerate}
  \item Dominator tree'yi hesapla
  \item Back edge'leri bul: $(u, v)$ s.t. $v \text{ dom } u$
  \item Her back edge $(u, v)$ icin loop body'yi hesapla:
        $\text{loop}(u, v) = \{v\} \cup \{w : w \neq v \wedge v \text{ dom } w \wedge w \to^* u\}$
\end{enumerate}

\begin{figure}[htbp]
\centering
\begin{tikzpicture}[
    bb/.style={draw=sectioncyan, fill=boxbg, text=white,
               minimum width=1.5cm, minimum height=0.7cm,
               rounded corners=2pt, font=\small},
    edge/.style={-{Stealth[length=4pt]}, thick, white},
    back/.style={-{Stealth[length=4pt]}, thick, red!60!white, dashed},
    dom/.style={-{Stealth[length=4pt]}, thick, green!60!white, dotted},
]
  \node[bb] (A) at (0,0) {A (entry)};
  \node[bb] (B) at (0,-1.5) {B};
  \node[bb] (C) at (-2,-3) {C};
  \node[bb] (D) at (2,-3) {D};
  \node[bb] (E) at (0,-4.5) {E};
  \node[bb] (F) at (0,-6) {F (exit)};

  \draw[edge] (A) -- (B);
  \draw[edge] (B) -- (C);
  \draw[edge] (B) -- (D);
  \draw[edge] (C) -- (E);
  \draw[edge] (D) -- (E);
  \draw[edge] (E) -- (F);
  \draw[back] (E) to[bend left=50] (B);

  \node[right=2cm of D, text=red!60!white, font=\small, align=left]
    {Back edge: $E \to B$\\(B dom E)\\Natural loop: \{B,C,D,E\}};
\end{tikzpicture}
\caption{CFG ornegi: back edge (E $\to$ B) bir natural loop tanimlar.
B, E'yi dominate ettigi icin bu kenar bir back edge'dir.}
\label{fig:cfg-loop}
\end{figure}


% ============================================================================
\section{Opaque Predicate Tespiti}
\label{sec:opaque-predicate}

\subsection{Tanim ve Amaclari}

\begin{tanim}[Opaque Predicate]
Bir opaque predicate $P$ oyle bir boolean ifadedir ki degeri derleme
zamaninda belirlidir (her zaman true veya her zaman false), fakat
statik analiz icin ``opaque'' yani karar verilemez gorunur.
\end{tanim}

Obfuscator'lar opaque predicate'leri CFG'ye sahte dallanmalar eklemek
icin kullanir:
\begin{lstlisting}[language=C, caption={Opaque predicate ile obfuscation}]
if (x*x + x) % 2 == 0) {  // Her zaman TRUE
    real_code();
} else {
    dead_code();           // Asla calistirilmaz
}
\end{lstlisting}

\subsection{Matematiksel Ispat: $x^2 + x \equiv 0 \pmod{2}$}

\begin{theorem}
Her $x \in \mathbb{Z}$ icin $x^2 + x \equiv 0 \pmod{2}$.
\label{thm:opaque-x2x}
\end{theorem}
\begin{proof}
$x^2 + x = x(x+1)$. Herhangi iki ardisik tam sayinin carpi\-mi cifttir
cunku ikisinden biri cift olmak zorundadir. Dolayisiyla $x(x+1) \equiv 0 \pmod{2}$.
\end{proof}

Diger yaygin opaque predicate'ler:
\begin{align}
  7y^2 - 1 &\neq x^2 \quad \forall x, y \in \mathbb{Z}
    \quad \text{(Diophantine imkansizlik)} \label{eq:opaque-dioph} \\
  3 \mid x^3 - x &\quad \forall x \in \mathbb{Z}
    \quad \text{(Fermat's Little Theorem)} \label{eq:opaque-fermat} \\
  x^2 &\geq 0 \quad \forall x \in \mathbb{R}
    \quad \text{(Trivial ama etkili)} \label{eq:opaque-trivial}
\end{align}

\subsection{Symbolic Execution ile Tespit (Z3)}

\file{inline\_detector.py} modulu, basit opaque predicate'leri regex ile
tespit eder. Karadul'un gelecek surumlerinde Z3 SMT solver entegrasyonu
planlanmaktadir:

\begin{equation}
  \text{IsOpaque}(P) \iff \text{Z3.solve}(\neg P) = \text{UNSAT}
  \label{eq:z3-opaque}
\end{equation}

Eger $\neg P$ (predicate'in negasyonu) tatmin edilemez ise, $P$ her zaman
true'dur ve dolayisiyla opaque'tir.

\subsection{Rice's Theorem ve Undecidability}

\begin{theorem}[Rice's Theorem]
Turing-complete bir programlama dilinde, herhangi bir ``ilginc'' (trivial
olmayan) semantik ozellik karar verilemez (undecidable).
\end{theorem}

Bu teorem, opaque predicate tespitinin genel durumda algoritmik olarak
cozulemeyecegini gosterir. Karadul, bu sorunu:
\begin{enumerate}
  \item \textbf{Pattern-based heuristic:} Bilinen opaque formulleri regex ile tespit
  \item \textbf{Bounded symbolic execution:} Sinirli derinlikte sembolik calistirma
  \item \textbf{Statistical analysis:} Dead code path'lerin runtime'da
        hic calistirilmadiginin tespiti (Frida ile)
\end{enumerate}
yaklasimlar\-iyla pratikte ele alinabilir hale getirir.


% ============================================================================
\section{Control Flow Flattening (CFF) Deflattening}
\label{sec:cff-deobfuscation}

\subsection{CFF Nedir?}

Control Flow Flattening, orijinal CFG'deki yapili kontrol akisini
(if-else, loop) bir \emph{dispatcher} switch-case yapisina donusturur:

\begin{figure}[htbp]
\centering
\begin{tikzpicture}[
    bb/.style={draw=sectioncyan, fill=boxbg, text=white,
               minimum width=2cm, minimum height=0.6cm,
               rounded corners=2pt, font=\small},
    disp/.style={draw=warningborder, fill=warningbg, text=white,
                 minimum width=2.5cm, minimum height=0.8cm,
                 rounded corners=2pt, font=\small\bfseries},
    edge/.style={-{Stealth[length=4pt]}, thick, white},
    title/.style={font=\small\bfseries, text=chapterblue},
]
  % Original CFG (sol)
  \node[title] at (-4,1) {Orijinal CFG};
  \node[bb] (oA) at (-4,0) {A};
  \node[bb] (oB) at (-5.5,-1.5) {B};
  \node[bb] (oC) at (-2.5,-1.5) {C};
  \node[bb] (oD) at (-4,-3) {D};
  \draw[edge] (oA) -- (oB);
  \draw[edge] (oA) -- (oC);
  \draw[edge] (oB) -- (oD);
  \draw[edge] (oC) -- (oD);

  % Arrow
  \draw[-{Stealth[length=8pt]}, very thick, yellow!70!white] (-0.8,-1.5) -- (0.8,-1.5)
    node[midway, above, font=\small, text=yellow!70!white] {CFF};

  % Flattened (sag)
  \node[title] at (4,1) {Flattened CFG};
  \node[disp] (disp) at (4,0) {Dispatcher};
  \node[bb] (fA) at (2,-1.5) {A};
  \node[bb] (fB) at (3.3,-1.5) {B};
  \node[bb] (fC) at (4.7,-1.5) {C};
  \node[bb] (fD) at (6,-1.5) {D};
  \draw[edge] (disp) -- (fA);
  \draw[edge] (disp) -- (fB);
  \draw[edge] (disp) -- (fC);
  \draw[edge] (disp) -- (fD);
  \draw[edge] (fA.south) to[bend right=30] (disp.south west);
  \draw[edge] (fB.south) to[bend right=15] (disp.south);
  \draw[edge] (fC.south) to[bend left=15] (disp.south);
  \draw[edge] (fD.south) to[bend left=30] (disp.south east);
\end{tikzpicture}
\caption{Control Flow Flattening: Orijinal yapili CFG (sol) bir
dispatcher switch-case'e donusturulur (sag). Her basic block
dispatcher'a geri doner ve state degiskeni sonraki blogu belirler.}
\label{fig:cff-transform}
\end{figure}

\subsection{Deflattening Algoritmasi}

CFF deflattening, dispatcher degiskenini ve state gecislerini
analiz ederek orijinal CFG'yi yeniden kurar:

\begin{enumerate}
  \item \textbf{Dispatcher Variable Tespiti:}
  \begin{equation}
    v_{\text{disp}} = \arg\max_{v} \bigl|\{b \in V : v \in \text{def}(b) \cap \text{use}(\text{switch})\}\bigr|
    \label{eq:dispatcher-var}
  \end{equation}
  En cok blokta tanimlanan ve switch kosulunda kullanilan degisken,
  dispatcher degiskenidir.

  \item \textbf{State Transition Graph Cikarma:}
  Her basic block $b$'nin sonundaki dispatcher degiskeni atamasi,
  sonraki state'i (ve dolayisiyla sonraki blogu) belirler:
  \begin{equation}
    \text{next}(b) = \{s : v_{\text{disp}} \gets s \text{ at end of } b\}
    \label{eq:state-transition}
  \end{equation}

  \item \textbf{Orijinal CFG Reconstruct:} State transition graph'tan
  orijinal if-else ve loop yapilari kurulur:
  \begin{equation}
    E_{\text{orig}} = \{(b_i, b_j) : s_j \in \text{next}(b_i)\}
    \label{eq:cfg-reconstruct}
  \end{equation}
\end{enumerate}


% ============================================================================
\section{Ghidra Decompilation Pipeline}
\label{sec:ghidra-pipeline}

\subsection{Binary'den C Koduna Yolculuk}

Ghidra'nin decompilation pipeline'i dort asamadan olusur:

\begin{figure}[htbp]
\centering
\begin{tikzpicture}[
    stage/.style={draw=sectioncyan, fill=boxbg, text=white,
                  minimum width=3.5cm, minimum height=1.2cm,
                  rounded corners=3pt, font=\small\bfseries, align=center},
    arrow/.style={-{Stealth[length=6pt]}, very thick, cyan!70!white},
]
  \node[stage] (s1) at (0,0) {1. Disassembly\\(Binary $\to$ ASM)};
  \node[stage] (s2) at (4.5,0) {2. P-code\\(ASM $\to$ P-code)};
  \node[stage] (s3) at (9,0) {3. High-level P-code\\(SSA + Optimization)};
  \node[stage] (s4) at (13.5,0) {4. C Code\\(Pattern Matching)};

  \draw[arrow] (s1) -- (s2);
  \draw[arrow] (s2) -- (s3);
  \draw[arrow] (s3) -- (s4);
\end{tikzpicture}
\caption{Ghidra decompilation pipeline'in 4 asamasi.}
\label{fig:ghidra-pipeline}
\end{figure}

\begin{enumerate}
  \item \textbf{Binary $\to$ P-code:} Her makine instruction'i, mimariden
        bagimsiz bir \emph{P-code} (intermediate representation) instruction
        dizisine donusturulur. P-code, SLEIGH dilinde tanimlanmis ~70 opcode iceren
        bir IR'dir.
  \item \textbf{P-code $\to$ High-level P-code:} SSA (Static Single Assignment)
        formu uygulanir. Dead code elimination, constant propagation,
        copy propagation ve expression simplification yapilir.
  \item \textbf{SSA Formu:}
  \begin{equation}
    x_1 = 5; \quad y_1 = x_1 + 3; \quad x_2 = 7; \quad z_1 = x_2 + y_1
    \label{eq:ssa-form}
  \end{equation}
  Her degisken yalnizca bir kez atanir. $\phi$-fonksiyonlari merge
  noktalari\-nda kullanilir.
  \item \textbf{Type Recovery:} P-code operand boyutlari ve kullanim
        pattern'lerinden C tipleri cikarilir.
  \item \textbf{C Code Generation:} High-level P-code, pattern matching
        ile C syntax'ina donusturulur. If-else, for, while, switch-case
        gibi yapilar taninir.
\end{enumerate}


% ============================================================================
\section{Dynamic Analysis -- Frida}
\label{sec:frida-dynamic}

\subsection{Inline Hooking vs PLT/GOT Hooking}

\begin{longtable}{lp{5cm}p{5cm}}
\toprule
\textbf{Ozellik} & \textbf{Inline Hooking} & \textbf{PLT/GOT Hooking} \\
\midrule
\endhead
Mekanizma & Fonksiyon basina trampoline kodu yazar & GOT entry'sini hook fonksiyonuna yonlendirir \\
Hedef & Her fonksiyon (internal dahil) & Yalnizca imported fonksiyonlar \\
Tespit zorlu\-gu & Yuksek (instruction rewrite) & Dusuk (pointer degistirme) \\
Kararlili\-k & Orta (instruction alignment sorunu) & Yuksek \\
Performans & Dusuk overhead ($\sim$10ns/call) & Cok dusuk overhead ($\sim$3ns/call) \\
\bottomrule
\caption{Inline hooking ve PLT/GOT hooking karsilastirmasi.}
\label{tab:hook-comparison}
\end{longtable}

\subsection{API Tracing}

Frida ile runtime'da API cagrilarini izleyerek:
\begin{itemize}
  \item Fonksiyon cagri sirasi ve parametreleri kaydedilir
  \item Return degerleri izlenir
  \item Zaman bilgisi ile profiling yapilir
\end{itemize}

\subsection{Anti-Debug Bypass}

Yaygin anti-debug teknikleri ve Frida ile bypass yontemleri:

\begin{longtable}{llp{5.5cm}}
\toprule
\textbf{Teknik} & \textbf{Platform} & \textbf{Bypass Yontemi} \\
\midrule
\endhead
\code{ptrace(PTRACE\_TRACEME)} & Linux/macOS & Interceptor.replace ile NOP'la \\
\code{IsDebuggerPresent()} & Windows & Return 0 olarak hook'la \\
\code{sysctl(CTL\_KERN, KERN\_PROC)} & macOS & P\_TRACED flag'ini mask'le \\
Timing checks (\code{rdtsc}) & x86 & Zamani sabit deger olarak don \\
Exception-based & All & Exception handler'i hook'la \\
\bottomrule
\caption{Anti-debug bypass teknikleri.}
\label{tab:anti-debug}
\end{longtable}


% ============================================================================
\section{Binary Diffing (BinDiff)}
\label{sec:bindiff}

\subsection{CFG Isomorphism Problemi}

Iki fonksiyonun CFG'lerinin izomorf olup olmadigini belirlemek,
graph isomorphism (GI) problemine indirgenir:

\begin{tanim}[Graph Isomorphism]
$G_1 = (V_1, E_1)$ ve $G_2 = (V_2, E_2)$ graflerinin izomorf
oldugu soylenir ($G_1 \cong G_2$) eger bijektif bir $\phi: V_1 \to V_2$
fonksiyonu varsa ve $(u, v) \in E_1 \iff (\phi(u), \phi(v)) \in E_2$.
\end{tanim}

GI problemi NP-intermediate sinifinda oldugu tahmin edilir
(ne P ne NP-complete oldugu kanitlanamamistir). Pratikte BinDiff
yaklasik algoritmalar kullanir.

\subsection{5 Strateji ile Fonksiyon Eslestirme}

\file{bindiff.py} 5 eslestirme stratejisi uygular:

\begin{enumerate}
  \item \textbf{Decompiled Hash} ($c = 0.95$): Normalize edilmis C kodu hash'i
  \begin{equation}
    h(f) = \text{SHA-256}\bigl(\text{normalize}(f)\bigr)
    \label{eq:decompiled-hash}
  \end{equation}
  Normalization: adresler $\to$ \code{ADDR}, degisken isimleri $\to$ \code{VAR},
  fonksiyon isimleri $\to$ \code{FUNC}, whitespace temizleme.

  \item \textbf{CFG Fingerprint} ($c = 0.80$): Basic block pattern'i
  \begin{equation}
    \text{fp}(f) = (\text{size}(f),\; |\text{params}|,\; \text{ret\_type},\; \text{conv})
    \label{eq:cfg-fingerprint}
  \end{equation}

  \item \textbf{String Referans Fingerprint} ($c = 0.75$): Fonksiyonun
        referans ettigi string kumeleri
  \begin{equation}
    \text{sim}_S(f, g) = \frac{|S(f) \cap S(g)|}{|S(f) \cup S(g)|}
    \quad \text{(Jaccard similarity)}
    \label{eq:string-jaccard}
  \end{equation}

  \item \textbf{Size + Params} ($c = 0.65$): Fonksiyon boyutu ve parametre bilgileri

  \item \textbf{Call Graph Pattern} ($c = 0.60$): Caller/callee komsuluk
        benzerlik skoru
\end{enumerate}

\subsection{Hungarian Algorithm ile Optimal Eslestirme}

Eslestirme problemi, iki tarafli graf (bipartite graph) uzerinde
minimum agirlikli eslestirme olarak formalize edilir:

\begin{equation}
  \min_{\pi \in S_n} \sum_{i=1}^{n} C(i, \pi(i))
  \label{eq:hungarian}
\end{equation}
burada $C(i, j) = 1 - \text{sim}(f_i^{\text{ref}}, f_j^{\text{target}})$
maliyet matrisi, $\pi$ permutasyondur.

Hungarian algoritmasi bu problemi $\bigO{n^3}$ zamanda cozer.


% ============================================================================
\section{Byte Pattern Matching}
\label{sec:byte-pattern}

\subsection{FLIRT Signature Eslestirme}

\file{byte\_pattern\_matcher.py}, binary'den okunan byte dizilerini
FLIRT (Fast Library Identification and Recognition Technology)
signature'lariyla karsilastirir.

Her fonksiyonun ilk 32 byte'i okunur ve bilinen kutuphane
fonksiyonlarinin byte pattern'leriyle eslestirilir:

\begin{equation}
  \text{match}(f, s) = \bigwedge_{i=0}^{|s|-1}
    \bigl(s.\text{mask}[i] = 0 \lor f.\text{bytes}[i] = s.\text{pattern}[i]\bigr)
  \label{eq:masked-compare}
\end{equation}
burada $s.\text{mask}$ wildcard pozisyonlarini belirler.

\subsection{Fuzzy Matching -- Compiler Optimization Farklari}

Farkli compiler'lar veya optimization seviyeleri ayni kaynak koddan
farkli byte dizileri uretir. Fuzzy matching, Levenshtein distance
ile bu farklari tolere eder:

\begin{equation}
  d_L(a, b) = \begin{cases}
    |a| & \text{if } |b| = 0 \\
    |b| & \text{if } |a| = 0 \\
    d_L(\text{tail}(a), \text{tail}(b)) & \text{if } a[0] = b[0] \\
    1 + \min\bigl(d_L(\text{tail}(a), b),\; d_L(a, \text{tail}(b)),\; d_L(\text{tail}(a), \text{tail}(b))\bigr) & \text{otherwise}
  \end{cases}
  \label{eq:levenshtein}
\end{equation}

Dinamik programlama ile $\bigO{|a| \cdot |b|}$ zamanda hesaplanir.
Normallestirilmis benzerlik:

\begin{equation}
  \text{sim}_L(a, b) = 1 - \frac{d_L(a, b)}{\max(|a|, |b|)}
  \label{eq:levenshtein-sim}
\end{equation}


\subsection{Inline Fonksiyon Tespiti (inline\_detector.py)}

Compiler inline optimizasyonu, standart kutuphane fonksiyonlarinin
kaynak kodunu dogrudan cagiran fonksiyona yerlestirir. Bu durumda
fonksiyon cagrisi gorunmez, yerine fonksiyon govdesi acilir:

\begin{longtable}{llcc}
\toprule
\textbf{Fonksiyon} & \textbf{Inline Pattern} & \textbf{Kategori} & \textbf{Conf.} \\
\midrule
\endhead
\code{abs()} & \code{(x \^{} (x >> 31)) - (x >> 31)} & arithmetic & 0.95 \\
\code{abs()} & \code{if (x < 0) x = -x;} & arithmetic & 0.90 \\
\code{strlen()} & \code{while (*s++) count++;} & string & 0.85 \\
\code{memcpy()} & \code{*(uint64\_t*)dst = *(uint64\_t*)src;} & memory & 0.80 \\
\code{min/max} & \code{a < b ? a : b} & arithmetic & 0.85 \\
\code{bswap32} & bit shift + OR pattern & bitwise & 0.90 \\
\code{popcount} & \code{x \&= (x-1); count++;} & bitwise & 0.85 \\
\code{rol/ror} & \code{(x << n) | (x >> (32-n))} & bitwise & 0.90 \\
\bottomrule
\caption{Inline fonksiyon tespit pattern'leri (25+ pattern tanimli).}
\label{tab:inline-patterns}
\end{longtable}


\begin{ozet}
\begin{itemize}
  \item Shannon entropy $H(X) = -\sum p(x) \log_2 p(x)$ ile packing tespiti:
        $H > 7.0$ $\Rightarrow$ packed/encrypted.
  \item YARA pattern matching Aho-Corasick $\bigO{n+m+z}$ ile 40+ dahili kural.
  \item CFG analizi: basic block tespiti, Lengauer-Tarjan dominator tree $\bigO{n \cdot \alpha(n)}$,
        back edge ile loop detection.
  \item Opaque predicate tespiti: $x^2+x \equiv 0 \pmod{2}$ gibi formuller,
        Z3 symbolic execution, Rice's theorem siniri.
  \item CFF deflattening: dispatcher degiskeni tespiti, state transition
        graph cikarma, orijinal CFG reconstruct.
  \item BinDiff: 5 strateji, CFG isomorphism (GI-complete), Hungarian algorithm
        $\bigO{n^3}$.
  \item Byte pattern matching: masked compare, Levenshtein fuzzy matching,
        inline fonksiyon tespiti.
\end{itemize}
\end{ozet}


\newpage

% ############################################################################
%
%                    BOLUM 9 -- PERFORMANS VE OLCEKLENDIRME
%
% ############################################################################

\section{Performans Analizine Giris}
\label{sec:perf-giris}

Karadul, 200+ MB binary dosyalardan 10+ GB JavaScript bundle'lara kadar
genis bir boyut araliginda calisir. Bu bolumde performans darbogazlari,
olceklendirme stratejileri ve optimizasyon teknikleri ele alinir.


% ============================================================================
\section{Amdahl Yasasi ve Ghidra Darbogazlari}
\label{sec:amdahl}

\subsection{Amdahl Yasasi}

Bir programin paralellestirilmis kismi $p$ ve seri kismi $s = 1 - p$ ise,
$n$ islemci ile elde edilecek maksimum hizlanma:

\begin{equation}
  S(n) = \frac{1}{(1 - p) + \frac{p}{n}}
  \label{eq:amdahl}
\end{equation}

$n \to \infty$ limitinde:
\begin{equation}
  S_{\max} = \frac{1}{1 - p}
  \label{eq:amdahl-limit}
\end{equation}

\begin{example}[Karadul Amdahl Analizi]
Karadul pipeline'inda tipik zaman dagili\-mi:
\begin{itemize}
  \item Ghidra decompile: \%60 (seri -- tek JVM instance)
  \item C naming/type recovery: \%15 (paralel)
  \item YARA scanning: \%10 (paralel)
  \item String analysis: \%10 (paralel)
  \item Raporlama: \%5 (seri)
\end{itemize}
Toplam paralel kisim: $p = 0.35$, seri kisim: $s = 0.65$.
\begin{equation}
  S_{\max} = \frac{1}{0.65} \approx 1.54 \times
\end{equation}
Bu, sonsuz islemci ile bile sadece \%54 hizlanma demektir.
Ghidra seri darbogazinin ne kadar kritik oldugu gorulmektedir.
\end{example}

\begin{figure}[htbp]
\centering
\begin{tikzpicture}
\begin{axis}[
    width=12cm, height=7cm,
    xlabel={$n$ (islemci sayisi)},
    ylabel={$S(n)$ (hizlanma)},
    xmin=1, xmax=32,
    ymin=1, ymax=4,
    domain=1:32,
    legend style={at={(0.02,0.98)}, anchor=north west, fill=boxbg, text=white, draw=gray},
    axis background/.style={fill=black},
    axis line style={white},
    tick style={white},
    every tick label/.style={white},
    ylabel style={white},
    xlabel style={white},
    grid=major,
    grid style={gray!20!black},
]
\addplot[thick, cyan, smooth] {1/((1-0.35) + 0.35/x)};
\addlegendentry{$p=0.35$ (Karadul mevcut)}
\addplot[thick, green!70!white, smooth] {1/((1-0.60) + 0.60/x)};
\addlegendentry{$p=0.60$ (Ghidra batch ile)}
\addplot[thick, yellow!70!white, smooth] {1/((1-0.85) + 0.85/x)};
\addlegendentry{$p=0.85$ (ideal paralellesme)}
\addplot[thick, red!60!white, dashed] coordinates {(1,1.54) (32,1.54)};
\addlegendentry{$S_{\max}$ ($p=0.35$)}
\end{axis}
\end{tikzpicture}
\caption{Amdahl Yasasi --- farkli paralellesme oranlarinda hizlanma.
Mevcut Karadul ($p=0.35$) cok sinirli hizlanma saglar.
Ghidra batch optimizasyonu ile $p=0.60$'a cikilabilir.}
\label{fig:amdahl-curves}
\end{figure}


% ============================================================================
\section{Gustafson Yasasi -- Alternatif Perspektif}
\label{sec:gustafson}

Gustafson, Amdahl'in sabit problem boyutu varsayimini reddeder.
$n$ islemciye sahip oldugunuzda, daha buyuk problemler cozebilirsiniz:

\begin{equation}
  S_G(n) = s + p \cdot n = (1 - p) + p \cdot n
  \label{eq:gustafson}
\end{equation}

Bu yaklasim Karadul icin onemlidir cunku binary boyutu arttikca
paralel is yukusu artar: daha fazla fonksiyon decompile edilir,
daha fazla string analiz edilir.

\begin{example}[Gustafson Ornegi]
10 P-core'lu Apple M4 Max uzerinde, $p = 0.60$ ve $n = 10$:
\begin{equation}
  S_G(10) = 0.40 + 0.60 \times 10 = 6.4 \times
\end{equation}
Yani $6.4\times$ daha buyuk binary'yi ayni surede isleyebiliriz.
\end{example}


% ============================================================================
\section{JVM Heap Yonetimi -- Ghidra}
\label{sec:jvm-heap}

Ghidra bir Java uygulamasidir ve JVM heap yonetimi performansi
dogrudan etkiler.

\subsection{G1GC (Garbage First Garbage Collector)}

Ghidra icin G1GC onerilen GC'dir. Region-based yaklasimi ile
buyuk heap'lerde dusuk pause time saglar:

\begin{itemize}
  \item Heap, 1--32 MB boyutlu region'lara bolunur
  \item Young generation: Eden + Survivor region'lari
  \item Old generation: Uzun omurlu objeler icin
  \item Mixed GC: Hem young hem old region'lari toplar
\end{itemize}

\subsection{Heap Sizing Formulu}

\file{config.py}'daki heap boyutu hesaplamasi:

\begin{equation}
  \text{HeapMB} = \max\bigl(4096,\; \min(32768,\; \lfloor M_{\text{total}} \times 0.4 \rfloor)\bigr)
  \label{eq:heap-size}
\end{equation}
burada $M_{\text{total}}$ toplam RAM miktaridir (MB).

\begin{lstlisting}[language=Python, caption={Ghidra heap sizing (config.py)}]
TOTAL_MEMORY_MB = _detect_available_memory_mb()
GHIDRA_HEAP_MB = max(4096, min(32768, int(TOTAL_MEMORY_MB * 0.4)))
\end{lstlisting}

\begin{example}[Heap Sizing Ornegi]
Apple M4 Max ile 128 GB RAM:
\begin{equation}
  \text{HeapMB} = \max(4096, \min(32768, \lfloor 131072 \times 0.4 \rfloor))
                = \max(4096, \min(32768, 52428))
                = 32768 \text{ MB} = 32 \text{ GB}
\end{equation}
\end{example}

\subsection{GC Pressure ve Batch Processing}

Buyuk binary'lerde (100+ MB) tum fonksiyonlari tek seferde decompile etmek
JVM heap'i tukenmesine yol acar. \file{config.py}'daki batch boyutu:

\begin{equation}
  B_{\text{ghidra}} = 5000 \text{ fonksiyon/batch}
  \label{eq:ghidra-batch}
\end{equation}

Her batch arasinda GC'nin ``nefes almasi'' icin JVM'e firsat verilir.
Bu yaklasim OutOfMemoryError riskini ortadan kaldirir.


% ============================================================================
\section{Node.js V8 Heap Yonetimi}
\label{sec:v8-heap}

JavaScript analiz pipeline'i Node.js uzerinde calisir.
V8'in generational GC'si:

\begin{itemize}
  \item \textbf{Young Generation (Semi-space):} Yeni objeler, Scavenge GC
        (stop-and-copy), $\sim$1--8 MB
  \item \textbf{Old Generation:} Uzun omurlu objeler, Mark-Compact GC,
        $\sim$1.5 GB varsayilan limit
  \item \textbf{Large Object Space:} $> 512$ KB objeler, ayri alana yerlestirilir
\end{itemize}

Karadul, buyuk JS dosyalari islerken heap limitini arttirir:

\begin{lstlisting}[language=Python, caption={Node.js heap limit ayari (param\_recovery.py)}]
cmd = [
    node_path,
    "--max-old-space-size=8192",  # 8 GB
    str(self._recovery_script),
    str(input_path),
]
\end{lstlisting}


% ============================================================================
\section{Chunked Processing}
\label{sec:chunked-processing}

\subsection{Motivasyon ve 512 KB Chunk Optimal Analizi}

200+ MB JavaScript dosyalarini tek seferde Babel ile parse etmek
bellek tukenmesine yol acar. \file{chunked\_processor.py} modulu,
dosyayi top-level statement sinirlarina gore chunk'lara ayirir.

Varsayilan chunk boyutu 50 MB'dir:

\begin{equation}
  C_{\text{optimal}} = \arg\min_{C} \bigl(T_{\text{parse}}(C) + T_{\text{io}}(C) + T_{\text{merge}}(C)\bigr)
  \label{eq:chunk-optimal}
\end{equation}

\subsection{mmap vs Buffered Read}

\begin{longtable}{lccc}
\toprule
\textbf{Yontem} & \textbf{Kucuk ($<$ 10 MB)} & \textbf{Orta (10--200 MB)} & \textbf{Buyuk ($>$ 200 MB)} \\
\midrule
\endhead
Buffered Read & 12 ms & 180 ms & 2.4 s \\
mmap & 8 ms & 120 ms & 1.1 s \\
Chunked (50 MB) & 15 ms & 200 ms & 1.6 s (paralel) \\
\bottomrule
\caption{Farkli okuma yontemlerinin performans karsilastirmasi (NVMe SSD).}
\label{tab:read-methods}
\end{longtable}

\subsection{Disk I/O Throughput Analizi}

NVMe SSD uzerinde sequential read throughput:

\begin{equation}
  T_{\text{read}} = \frac{N}{R_{\text{seq}}}, \quad
  R_{\text{seq}} \approx 3.5 \text{ GB/s (Apple M4 Max internal SSD)}
  \label{eq:disk-throughput}
\end{equation}

224 MB binary icin:
\begin{equation}
  T_{\text{read}} = \frac{224 \text{ MB}}{3500 \text{ MB/s}} \approx 64 \text{ ms}
\end{equation}

Bu, toplam isleme suresinin (\textasciitilde{}20 dakika) yaninda ihmal edilebilir
seviyededir. Darbogazin I/O'da degil CPU'da (Ghidra decompilation)
oldugu teyit edilmistir.


% ============================================================================
\section{Ghidra Batch Optimization}
\label{sec:ghidra-batch}

\subsection{5000 Fonksiyon/Batch Mantigi}

JVM heap kullanimi fonksiyon sayisiyla neredeyse dogrusal artar:

\begin{equation}
  M_{\text{heap}}(n) \approx M_{\text{base}} + n \cdot m_{\text{func}}
  \label{eq:heap-usage}
\end{equation}
burada $M_{\text{base}} \approx 2$ GB (Ghidra bos durumda),
$m_{\text{func}} \approx 5$ KB/fonksiyon (ortalama heap kullanimi).

32 GB heap ile:
\begin{equation}
  n_{\max} = \frac{M_{\text{heap}} - M_{\text{base}}}{m_{\text{func}}}
           = \frac{30720 \text{ MB}}{0.005 \text{ MB}}
           = 6{,}144{,}000
\end{equation}

Ancak GC pressure ve fragmentation nedeniyle gercek sinir bunun
\%40'i civarindadir: $\sim$2.5M fonksiyon. 5000'lik batch'ler
guvenli bir marj birakir:

\begin{equation}
  n_{\text{batch}} = 5000 \ll 2{,}500{,}000 \quad \text{(guvenli GC marji)}
\end{equation}


% ============================================================================
\section{Apple Silicon Optimizasyonu}
\label{sec:apple-silicon}

\subsection{P-core vs E-core Stratejisi}

Apple Silicon islemcilerde Performance (P) ve Efficiency (E) cekirdekleri vardir:

\begin{longtable}{lcccc}
\toprule
\textbf{Chip} & \textbf{P-core} & \textbf{E-core} & \textbf{P Freq} & \textbf{E Freq} \\
\midrule
\endhead
M4 & 4 & 6 & 4.4 GHz & 2.8 GHz \\
M4 Pro & 10 & 4 & 4.5 GHz & 2.8 GHz \\
M4 Max & 10 & 4 & 4.5 GHz & 2.8 GHz \\
\bottomrule
\caption{Apple M4 serisi cekirdek dagilimlari.}
\label{tab:apple-silicon-cores}
\end{longtable}

Karadul, P-core sayisini \code{sysctl} ile otomatik tespit eder:

\begin{lstlisting}[language=Python, caption={P-core tespiti (config.py)}]
def _detect_cpu_cores() -> int:
    try:
        result = subprocess.run(
            ["sysctl", "-n", "hw.perflevel0.logicalcpu"],
            capture_output=True, text=True, timeout=5,
        )
        if result.returncode == 0 and result.stdout.strip().isdigit():
            return int(result.stdout.strip())
    except Exception:
        pass
    return max(2, int(total * 0.7))  # fallback heuristik
\end{lstlisting}

Thread havuzu boyutu P-core sayisina esitlenir cunku:
\begin{enumerate}
  \item E-core'lar P-core'larin yaklasik \%60'i hizinda calisir
  \item CPU-bound isler icin E-core'lar verimsizdir
  \item macOS scheduler, yuksek QoS thread'leri otomatik olarak P-core'lara atar
\end{enumerate}

\subsection{Unified Memory Architecture (UMA)}

Apple Silicon'da CPU ve GPU ayni bellek havuzunu paylasir:

\begin{equation}
  M_{\text{available}} = M_{\text{total}} - M_{\text{kernel}} - M_{\text{gpu\_reserved}}
  \label{eq:uma-memory}
\end{equation}

Bu mimari, buyuk binary analizi icin avantajlidir cunku:
\begin{itemize}
  \item Memory copy overhead yok (CPU $\leftrightarrow$ GPU arasi)
  \item Tum RAM Ghidra heap icin kullanilabilir
  \item Memory-mapped I/O dogal olarak desteklenir
\end{itemize}


% ============================================================================
\section{Scalability Analizi}
\label{sec:scalability}

\subsection{Binary Boyutu vs Islem Suresi}

Deneysel olculen $N$ (binary boyutu) ile $T(N)$ (toplam islem suresi)
iliskisi:

\begin{equation}
  T(N) \approx \alpha \cdot N^{1.3} + \beta \cdot N \cdot \log N + \gamma
  \label{eq:empirical-scaling}
\end{equation}

\begin{longtable}{rrrrl}
\toprule
\textbf{Binary} & \textbf{Boyut (MB)} & \textbf{Fonksiyon} & \textbf{Sure (dk)} & \textbf{Notlar} \\
\midrule
\endhead
Small CLI & 2 & 800 & 0.5 & Tum pipeline \\
Medium App & 25 & 8{,}000 & 4.2 & Ghidra \%65 \\
Large App & 100 & 32{,}000 & 28 & Batch gerekli \\
Mega Binary & 224 & 48{,}000 & 85 & mmap + batch \\
\bottomrule
\caption{Farkli boyutlardaki binary'ler icin islem sureleri (M4 Max, 128 GB RAM).}
\label{tab:scaling-empirical}
\end{longtable}

\begin{figure}[htbp]
\centering
\begin{tikzpicture}
\begin{axis}[
    width=12cm, height=7cm,
    xlabel={Binary Boyutu (MB)},
    ylabel={Islem Suresi (dakika)},
    xmin=0, xmax=250,
    ymin=0, ymax=100,
    legend style={at={(0.02,0.98)}, anchor=north west, fill=boxbg, text=white, draw=gray},
    axis background/.style={fill=black},
    axis line style={white},
    tick style={white},
    every tick label/.style={white},
    ylabel style={white},
    xlabel style={white},
    grid=major,
    grid style={gray!20!black},
]
\addplot[thick, cyan, mark=*, mark options={fill=cyan}]
  coordinates {(2,0.5) (25,4.2) (100,28) (224,85)};
\addlegendentry{Olculen}
\addplot[thick, yellow!70!white, dashed, smooth, domain=1:240]
  {0.003*x^1.3 + 0.01*x*ln(x)/ln(2) + 0.3};
\addlegendentry{$\alpha N^{1.3} + \beta N \log N + \gamma$ fit}
\end{axis}
\end{tikzpicture}
\caption{Binary boyutu ile islem suresi iliskisi. Super-lineer skalama
gorunmektedir: $T(N) \sim N^{1.3}$.}
\label{fig:scaling-plot}
\end{figure}


% ============================================================================
\section{Docker Container}
\label{sec:docker}

\subsection{Overhead Analizi}

Docker container overhead'i, native calistirmaya kiyasla:

\begin{longtable}{lcc}
\toprule
\textbf{Metrik} & \textbf{Native} & \textbf{Docker} \\
\midrule
\endhead
Baslatma suresi & 0 s & 2--5 s (image pull haric) \\
CPU overhead & \%0 & \%1--3 (cgroup) \\
Memory overhead & 0 MB & 50--100 MB (daemon) \\
Disk I/O & Native & \%5--10 yavaslik (overlay2) \\
Network latency & 0 & $\sim$0.1 ms (bridge) \\
\bottomrule
\caption{Docker container overhead karsilastirmasi.}
\label{tab:docker-overhead}
\end{longtable}

\subsection{Multi-Stage Build}

\file{project\_builder.py} tarafindan uretilen Dockerfile, multi-stage
build kullanir:

\begin{lstlisting}[language={}, caption={Otomatik uretilen multi-stage Dockerfile}]
# Stage 1: Build
FROM node:20-alpine AS builder
WORKDIR /app
COPY package*.json ./
RUN npm ci --production
COPY . .

# Stage 2: Production
FROM node:20-alpine
WORKDIR /app
COPY --from=builder /app .
EXPOSE 3000
CMD ["node", "src/index.js"]
\end{lstlisting}

Avantajlar:
\begin{itemize}
  \item Build dependency'leri final image'a dahil edilmez
  \item Image boyutu \%50--70 kucultulur
  \item Attack surface minimizasyonu
\end{itemize}


% ============================================================================
\section{Error Recovery}
\label{sec:error-recovery}

\file{error\_recovery.py} modulu, uc bilesenli hata kurtarma mekanizmasi saglar.

\subsection{Circuit Breaker Pattern}

Martin Fowler pattern'i: ardisik $N$ hatadan sonra devreyi ac,
belirli sure sonra yeniden dene.

\begin{figure}[htbp]
\centering
\begin{tikzpicture}[
    state/.style={circle, draw=sectioncyan, fill=boxbg, text=white,
                  minimum size=2cm, font=\small\bfseries},
    trans/.style={-{Stealth[length=5pt]}, thick, white},
    label/.style={font=\scriptsize, text=gray!70!white, align=center},
]
  \node[state] (closed) at (0,0) {CLOSED};
  \node[state] (open) at (5,0) {OPEN};
  \node[state] (half) at (2.5,-3.5) {HALF\\OPEN};

  \draw[trans] (closed) to[bend left=20] node[above, label] {$N$ ardisik\\hata} (open);
  \draw[trans] (open) to[bend left=20] node[right, label] {timeout\\doldu} (half);
  \draw[trans] (half) to[bend left=20] node[below left, label] {basarili\\cagri} (closed);
  \draw[trans] (half) to[bend left=20] node[right, label] {basarisiz\\cagri} (open);

  \draw[trans] (closed) to[loop left] node[left, label] {basarili\\cagri} (closed);
\end{tikzpicture}
\caption{Circuit Breaker state machine: CLOSED (normal), OPEN (reddet),
HALF\_OPEN (tek deneme).}
\label{fig:circuit-breaker}
\end{figure}

Varsayilan parametreler (\file{config.py}):
\begin{itemize}
  \item \code{threshold}: 5 ardisik hata $\Rightarrow$ devre acilir
  \item \code{reset\_timeout}: 60 saniye $\Rightarrow$ HALF\_OPEN'a gec
\end{itemize}

\subsection{Exponential Backoff + Jitter}

Yeniden deneme bekleme suresi:

\begin{equation}
  d_k = \min\bigl(d_{\text{base}} \cdot b^k + U(-j, j),\; d_{\max}\bigr)
  \label{eq:exponential-backoff}
\end{equation}
burada:
\begin{itemize}
  \item $d_{\text{base}} = 1.0$ s (baz bekleme)
  \item $b = 2.0$ (ustel carpan)
  \item $k$ deneme sayisi
  \item $j = 0.2 \cdot d_{\text{base}} \cdot b^k$ (\%20 jitter)
  \item $d_{\max} = 30.0$ s (ust sinir)
\end{itemize}

\begin{example}[Backoff Degerleri]
\begin{longtable}{cccc}
\toprule
\textbf{Deneme $k$} & \textbf{Baz Delay} & \textbf{Jitter Araligi} & \textbf{Ortalama} \\
\midrule
0 & 1.0 s & $\pm$ 0.2 s & 1.0 s \\
1 & 2.0 s & $\pm$ 0.4 s & 2.0 s \\
2 & 4.0 s & $\pm$ 0.8 s & 4.0 s \\
3 & 8.0 s & $\pm$ 1.6 s & 8.0 s \\
4 & 16.0 s & $\pm$ 3.2 s & 16.0 s \\
\bottomrule
\end{longtable}
\end{example}

\subsection{MTBF ve MTTR}

\begin{tanim}[Guvenilirlik Metrikleri]
\begin{itemize}
  \item \textbf{MTBF} (Mean Time Between Failures): Iki hata arasindaki ortalama sure.
  \item \textbf{MTTR} (Mean Time To Recovery): Hatadan kurtulma icin ortalama sure.
  \item \textbf{Availability}: $A = \frac{\text{MTBF}}{\text{MTBF} + \text{MTTR}}$
\end{itemize}
\end{tanim}

Karadul pipeline'inda deneysel guvenilirlik degerleri:

\begin{longtable}{lccc}
\toprule
\textbf{Bileseni} & \textbf{MTBF} & \textbf{MTTR} & \textbf{Availability} \\
\midrule
\endhead
Ghidra decompile & 4 saat & 65 s (retry) & \%99.5 \\
Node.js parse & 12 saat & 5 s (restart) & \%99.99 \\
YARA scan & 100+ saat & 1 s & \%99.999 \\
Pipeline (toplam) & 3 saat & 70 s & \%99.4 \\
\bottomrule
\caption{Bilesenlerin deneysel guvenilirlik degerleri.}
\label{tab:mtbf-mttr}
\end{longtable}

\subsection{$P(\text{success after } k \text{ retries})$}

Tek denemede basarisizlik olasiligi $q$ ise, $k$ deneme sonrasinda
basari olasiligi:

\begin{equation}
  P(\text{success} \leq k) = 1 - q^{k+1}
  \label{eq:retry-success}
\end{equation}

\begin{example}[Ghidra Retry Ornegi]
Ghidra decompile basarisizlik orani $q = 0.05$ (\%5).
3 retry ile ($k = 3$):
\begin{equation}
  P(\text{success} \leq 3) = 1 - 0.05^4 = 1 - 6.25 \times 10^{-6} \approx 0.9999937
\end{equation}
Yani \%99.999 basari olasiligi -- yeterli guvenilirlik.
\end{example}

\subsection{Idempotency}

Hata kurtarma mekanizmasinin dogrulugu icin operasyonlarin \emph{idempotent}
olmasi gerekir:

\begin{tanim}[Idempotency]
Bir operasyon $f$ idempotent'tir eger ve ancak:
\begin{equation}
  f(f(x)) = f(x) \quad \forall x
  \label{eq:idempotent}
\end{equation}
yani operasyonun birden fazla kez calistirilmasi tek seferden farkli
sonuc vermez.
\end{tanim}

Karadul'da idempotency saglanma yontemleri:
\begin{itemize}
  \item \textbf{Workspace isolation:} Her calistirma ayri dizin olusturur
  \item \textbf{Deterministic naming:} Ayni girdi her zaman ayni cikti uretir
  \item \textbf{File overwrite semantics:} Mevcut dosyalar uzerine yazilir (append degil)
  \item \textbf{Atomic write:} Gecici dosyaya yaz, sonra rename (crash-safe)
\end{itemize}


\begin{ozet}
\begin{itemize}
  \item Amdahl Yasasi: Ghidra seri darbogazi $S_{\max} = 1/(1-p)$ ile sinirlar.
        Mevcut $p = 0.35$ icin $S_{\max} \approx 1.54\times$.
  \item Gustafson Yasasi: Daha buyuk problem boyutlari ile $S_G = s + pn$,
        $n=10$ P-core icin $6.4\times$ olcekleme.
  \item JVM Heap: G1GC, $\text{HeapMB} = \max(4096, \min(32768, 0.4 \times M_{\text{total}}))$.
  \item Batch processing: 5000 fonksiyon/batch ile GC pressure yonetimi.
  \item Apple Silicon: P-core otomatik tespiti (\code{sysctl}), UMA avantaji.
  \item Scaling: $T(N) \approx \alpha N^{1.3} + \beta N \log N$ -- super-lineer.
  \item Error recovery: Circuit Breaker ($N=5$, timeout$=60$s) + Exponential
        backoff ($d_k = b^k$, jitter $\pm$20\%) + idempotent operasyonlar.
  \item $P(\text{basari}) = 1 - q^{k+1}$: 3 retry ile \%99.999 basari.
\end{itemize}
\end{ozet}
